% Generated from article.tex and sections/*.tex by build.py.
\documentclass[11pt,a4paper]{article}
\usepackage[T1]{fontenc}
\usepackage[utf8]{inputenc}
\usepackage{lmodern}
\usepackage[margin=26mm,headheight=15pt]{geometry}
\usepackage{amsmath,amssymb,amsthm,mathtools,mathrsfs}
\usepackage{booktabs,longtable,array,enumitem,microtype,xcolor,fancyhdr,xurl}
\usepackage[colorlinks=true,linkcolor=blue!45!black,citecolor=blue!45!black,urlcolor=blue!45!black]{hyperref}
\usepackage{bookmark}
\numberwithin{equation}{section}
\newtheorem{theorem}{Theorem}[section]
\newtheorem{proposition}[theorem]{Proposition}
\newtheorem{lemma}[theorem]{Lemma}
\newtheorem{corollary}[theorem]{Corollary}
\theoremstyle{definition}
\newtheorem{definition}[theorem]{Definition}
\newtheorem{conjecture}[theorem]{Conjecture}
\theoremstyle{remark}
\newtheorem{remark}[theorem]{Remark}
\newtheorem{example}[theorem]{Example}
\DeclareMathOperator{\FP}{FP}
\DeclareMathOperator{\CT}{CT}
\DeclareMathOperator*{\Res}{Res}
\DeclareMathOperator{\Li}{Li}
\DeclareMathOperator{\Log}{Log}
\DeclareMathOperator{\re}{Re}
\DeclareMathOperator{\im}{Im}
\DeclareMathOperator{\rank}{rank}
\DeclareMathOperator{\wt}{wt}
\newcommand{\dd}{\,\mathrm d}
\newcommand{\C}{\mathbb C}
\newcommand{\R}{\mathbb R}
\newcommand{\Q}{\mathbb Q}
\newcommand{\N}{\mathbb N}
\newcommand{\Z}{\mathbb Z}
\newcommand{\e}{\mathrm e}
\newcommand{\ii}{\mathrm i}
\newcommand{\cI}{\mathcal I}
\newcommand{\cD}{\mathcal D}
\newcommand{\cB}{\mathcal B}
\newcommand{\sourcecommit}{f7e35be357654198c1546fd40f363694c5bcc2cc}
\newcommand{\src}[1]{\nolinkurl{#1}}
\setlength{\emergencystretch}{3em}
\setlength{\parskip}{0.25em}
\setlist{itemsep=3pt,topsep=5pt}
\allowdisplaybreaks[2]
\pagestyle{fancy}
\fancyhf{}
\fancyhead[L]{\small Differential Polygamma and Lerch Identities}
\fancyhead[R]{\small ProveIt research continuation}
\fancyfoot[C]{\thepage}
\hypersetup{pdftitle={Differential Polygamma and Lerch Identities},pdfauthor={Research continuation prepared with ChatGPT for Vladimir Reshetnikov and ProveIt}}
\begin{document}
\hypersetup{pageanchor=false}
\begin{titlepage}
\vspace*{10mm}
\noindent{\large\sffamily PROVEIT / RESEARCH CONTINUATION}
\vspace{17mm}

\noindent{\Huge\bfseries Differential Polygamma\\[3mm]
and Lerch Identities}
\vspace{8mm}

\noindent{\Large Exact residue cancellation, higher Gamma moments,\\[2mm]
spectral derivatives, and harmonic-tail reconstruction}
\vspace{14mm}

\noindent Prepared with ChatGPT for Vladimir Reshetnikov and the ProveIt project\\
11 October 2026
\vspace{11mm}

\begin{abstract}
We continue six newly submitted reports and the current ProveIt polylogarithm
manuscript. The main theorem extends the normalized digamma-power calculus
to every differential polynomial in the digamma function. Its asymptotic
residue and its large-circle constant depend only on the scalar digamma
part. In particular, every monomial containing a positive-order polygamma
has an exact finite-part integral expressed by convergent multiple-zeta
values and convergent spectral derivatives, with no singular Laurent
coordinate. A triangular integration-by-parts theorem gives a canonical
formal basis and retains every endpoint constant. We give normalized
primitives, a common-shift endpoint law, polylogarithmic kernels, and an
explicit formula for every bilinear polygamma integral.

The Gamma--Appell continuation isolates the highest height-one Laurent
coordinate at every degree and makes the remaining coordinates convergent;
the sixth and seventh degrees are worked out. A second continuation
evaluates all spectral jets of the centered continuous-depth Lerch family
on its odd lattice, including mixed depth derivatives and Stieltjes
specializations. Finally, an exact determinant formula proves the
optimal-prefix conjecture for every neighboring pair of even harmonic-tail
orders. The unrestricted unequal-order conjecture and the Gaussian
$S_6,S_8$ candidates remain open. Proofs, source comparisons, and independent
symbolic and numerical checks are supplied; numerical agreement is never
used as a proof of an identity.
\end{abstract}
\vfill
\noindent\small Source baseline: ProveIt revision
\texttt{\sourcecommit}.\\
The contribution is measured against the inspected sources. A claim of
historical priority for every displayed specialization is not made.
\end{titlepage}
\hypersetup{pageanchor=true}
\pagenumbering{roman}
\tableofcontents
\newpage
\pagenumbering{arabic}

% BEGIN sections/01_scope.tex
\section{Research targets and conventions}
\label{poly:scope}

The newest incoming reports make the next exact questions unusually
specific. The higher-digamma report asks for a calculus of arbitrary
differential polynomials, for common shifted endpoints, and for explicit
sixth and seventh moments \cite{Higher}. The canonical fifth-moment report
asks which differential monomials survive integration by parts, with the
finite endpoint value retained, and conjectures that an optimal initial
set of harmonic-tail moments has full formal rank \cite{Fifth}.
The continuous-depth report asks for spectral derivatives at its odd
evaluation lattice \cite{Continuous}. These are the principal targets here.

The normalized Mittag--Leffler expansion for scalar powers, the
height-one Gamma integral, the centered Lerch reflection, and the original
two-tail moment formulas are antecedents. We reproduce the parts needed
to fix conventions and then prove their extensions. This distinction is
important because the repository already contains several independently
derived versions of related calculations \cite{Repo,Mixed,Tilted,Algebraic}.
The present audit is targeted, not an audit of every analytic assertion in
the collection.

\begin{center}
\small
\begin{tabular}{p{0.27\textwidth}p{0.62\textwidth}}
\toprule
Source question & Result and remaining scope\\
\midrule
Differential-polynomial moments & A complete residue and finite-part
theorem, with a finite reduction to convergent zeta coordinates for every
mixed monomial.\\[4pt]
Integration by parts at all weights & An explicit direct-sum normal form;
its dimension is a formal algebraic dimension, not independence of
the resulting numbers.\\[4pt]
Higher scalar Gamma normalization & Only the highest height-one Laurent
coordinate is singular; all remaining coordinates have convergent
representatives. Arithmetic reduction of those coordinates remains a
separate question.\\[4pt]
Odd-lattice Lerch spectral jets & Every order and mixed depth derivative,
with explicit local Gamma coefficients and convergent spectral evaluators.\\[4pt]
Optimal initial tail moments & Proved for all neighboring even orders
$p=q-2$. Arbitrary even gaps remain open.\\
\bottomrule
\end{tabular}
\end{center}

\subsection{Functions, finite parts, and spectral conventions}

Write $\psi=\Gamma'/\Gamma$ and $\psi^{(r)}=\dd^r\psi/\dd x^r$,
with $\psi^{(0)}=\psi$. Let
\[
 H_n^{(r)}=\sum_{k=1}^n k^{-r},\qquad H_n=H_n^{(1)},\qquad
 A_r(n)=\zeta(r)-H_n^{(r)}\quad(r\ge2).
\]
The empty sum is zero. A strict word has the convention
\[
 H_n(w_1,\ldots,w_d)=
 \sum_{n\ge k_1>\cdots>k_d\ge1}\prod_{j=1}^d k_j^{-w_j},
 \qquad H_n(\varnothing)=1.
\]
For convergent arguments, $Z(s_1,\ldots,s_d)$ is the corresponding
infinite strict multiple-zeta series. A derivative $\partial_j Z$
means differentiation in its $j$th spectral argument; its summand contains
$-\log k_j$. Classical definitions and normalizations of Gamma, Hurwitz
zeta, and Lerch's transcendent are those of the DLMF
\cite{DLMFGamma,DLMFZeta,DLMFLerch}.

The Stieltjes convention is
\begin{equation}\label{poly:stieltjes}
 \zeta(1+s,a)=\frac1s+
 \sum_{j\ge0}\frac{(-1)^j}{j!}\gamma_j(a)s^j,
 \qquad \re a>0.
\end{equation}
We write $\gamma_j=\gamma_j(1)$ and $\gamma=\gamma_0$.
The notation $[s]F(s)$ denotes a Laurent coefficient. It agrees with
$F'(0)$ only when $F$ is holomorphic at zero.

Suppose $f(x)=\sum_{j=1}^M p_jx^{-j}+O(1)$ near $0$. Our finite part is
\begin{equation}\label{poly:fp}
 \FP\int_0^1f(x)\dd x=
 \lim_{\varepsilon\downarrow0}
 \left\{\int_\varepsilon^1f(x)\dd x
 -\sum_{j=2}^M\frac{p_j\varepsilon^{1-j}}{j-1}
 +p_1\log\varepsilon\right\}.
\end{equation}
The coordinate is exactly $x$ and the upper endpoint exactly $1$.
Every finite-part statement below uses this definition. Normalized
primitives use principal logarithms on $\C\setminus(-\infty,0]$.

For a polynomial $P\in\C[u_0,u_1,\ldots,u_R]$, set
\begin{align*}
 f_P(x)&=P(\psi(x),\psi'(x),\ldots,\psi^{(R)}(x)),&
 \phi_P(Y)&=P(Y,0,\ldots,0),\\
 \cI(P)&=\FP\int_0^1f_P(x)\dd x.
\end{align*}
We call $P$ mixed when $\phi_P=0$. This is a precise algebraic condition:
each of its monomials contains at least one $u_r$ with $r\ge1$.
The word does not assert that the associated numerical periods are
independent. Give $u_r$ differential weight $r+1$.

% END sections/01_scope.tex


% BEGIN sections/02_differential_calculus.tex
\section{A differential-polynomial finite-part calculus}
\label{poly:calculus}

\subsection{Finite local data and the global normalization}

For $n\ge0$, put $b_n=H_n-\gamma$ and
$c_{r,n}=H_n^{(r)}+(-1)^r\zeta(r)$ for $r\ge2$. The familiar
Gamma recurrence and Taylor expansion give
\begin{equation}\label{poly:localpsi}
 v_n(\varepsilon):=\psi(-n+\varepsilon)
 =-\varepsilon^{-1}+b_n+
      \sum_{r\ge2}c_{r,n}\varepsilon^{r-1}.
\end{equation}
This series converges for $0<|\varepsilon|<1$; its differentiated
series give all polygammas. Let $M$ be an upper bound for the differential
weights of the monomials of $P$. Define
\begin{equation}\label{poly:poles}
 p_j(n)=[\varepsilon^{-j}]P(v_n,v_n',\ldots,v_n^{(R)}),\quad
 P_n(z)=\sum_{j=1}^M\frac{p_j(n)}{(z+n)^j},\quad
 \kappa_P=[\varepsilon^0]f_P(\varepsilon).
\end{equation}
All coefficients are finite polynomials in the stated harmonic numbers
and zeta values. To find the coefficients through the constant term,
one may truncate $v_n$ after $\varepsilon^{M-1}$ before differentiating.
For a homogeneous monomial of weight $M$, the regular coefficient of
harmonic order $r$ raises the exponent relative to its all-pole term
by $r$; hence orders above $M$ cannot contribute. Linearity handles
inhomogeneous polynomials.

\begin{theorem}[Normalized differential-polynomial expansion]
\label{poly:ML}
Away from the nonpositive integers,
\begin{equation}\label{poly:MLeq}
 f_P(z)=P_0(z)+\kappa_P+
              \sum_{n\ge1}\bigl(P_n(z)-P_n(0)\bigr).
\end{equation}
The series converges normally on compact subsets of its domain.
The entire remainder is exactly the displayed constant.
\end{theorem}

\begin{proof}
Each $p_j(n)$ grows at most as a fixed power of $1+\log(n+1)$:
only $H_n$ is unbounded, whereas the higher harmonic numbers are bounded.
On a fixed compact set,
$(n+z)^{-j}-n^{-j}=O(n^{-j-1})$. This proves normal convergence.
To determine the remainder, put $g(w)=f_P(w)-P_0(w)$ and integrate
$zg(w)/(w(w-z))$ on $|w|=N+1/2$.

The right-half-plane asymptotics give $\psi(w)=O(1+\log N)$ and
bounded fixed positive derivatives. On the left use
$\psi(w)=\psi(1-w)-\pi\cot(\pi w)$ and its derivatives.
The cotangent and each fixed derivative are uniformly bounded on these
circles. For $|\im w|\ge1$ this follows from their exponential
representations; for $|\im w|<1$ the circles remain a fixed positive
distance from the integers. Thus $g=O((1+\log N)^d)$ for some fixed
$d$, and the contour integral tends to zero. The residues at $z$ and
$0$ are $g(z)$ and $-\kappa_P$; a pole at $-n$ contributes
$-P_n(z)+P_n(0)$. The residue theorem gives the assertion.
\end{proof}

\subsection{The scalar projection controls both boundary obstructions}

Define polynomials in $L$ by
\begin{align}
 C_P(L)&=\frac1{2\pi}\int_{-\pi}^{\pi}\phi_P(L+\ii\theta)\dd\theta,
 & C_P&=C_P(0),\label{poly:C}\\
 R_P(L)&=-\frac{\phi_P(L+\ii\pi)-\phi_P(L-\ii\pi)}{2\ii\pi}.
 \label{poly:R}
\end{align}
The integrals are integrals of polynomials, so these are finite algebraic
formulas. If $\phi_P(Y)=\sum a_kY^k$, then
\begin{equation}\label{poly:Cexplicit}
 C_P=\sum_{m\ge0}a_{2m}\frac{(-1)^m\pi^{2m}}{2m+1}.
\end{equation}

\begin{theorem}[Scalar projection and residue cancellation]
\label{poly:filter}
For every fixed $P$, with $r_P(n)=p_1(n)$, there is an integer $d$ such that
\begin{align}
 \kappa_P-\sum_{n=1}^NP_n(0)&=C_P(\log N)+o(1),\label{poly:circle}\\
 r_P(n)&=R_P(\log n)+O((1+\log n)^d/n),\label{poly:rasym}\\
 r_P(n)-r_P(n-1)&=\frac{R_P'(\log n)}n
                      +O((1+\log n)^d/n^2).\label{poly:deltarasym}
\end{align}
In particular, a mixed polynomial has $C_P=R_P=0$.
\end{theorem}

\begin{proof}
The first identity follows by integrating $f_P(w)/w$ on the same
circles. Away from an arc of angular length $O(\log N/N)$ around the
negative real axis, reflection and the standard sectorial asymptotics
give $\psi(w)=\Log w+O(1/N)$ and
$\psi^{(r)}(w)=O(N^{-r})+O(N^{-2\pi})$ for fixed $r\ge1$.
On the exceptional arc all fixed polygammas are bounded and the digamma
is $O(\log N)$. Every mixed monomial consequently has angular average
$o(1)$. The scalar terms have the stated logarithmic average. Replacing
$\log(N+1/2)$ by $\log N$ changes the answer by $o(1)$.

For the residue, take the constant-tail specialization of the local
series: $b_n$ tends formally to $L$, and
$c_{r,n}$ tends to $(1+(-1)^r)\zeta(r)$. The resulting series is
$L+y$, where $y=-\pi\cot(\pi\varepsilon)$. Its derivatives satisfy
\[
 y'=y^2+\pi^2,
 \qquad y^{(r)}=Q_r(y),\quad
 Q_1(Y)=Y^2+\pi^2,\quad
 Q_{r+1}(Y)=(Y^2+\pi^2)Q_r'(Y).
\]
Every $Q_r$ with $r\ge1$ is divisible by $Y^2+\pi^2$.
Under the local change of variable $Y=y$, residue extraction is
\begin{equation}\label{poly:reschange}
 \Res_{\varepsilon=0}P(L+y,y',\ldots)
 =-[Y^{-1}]
 \frac{P(L+Y,Q_1(Y),Q_2(Y),\ldots)}{Y^2+\pi^2},
\end{equation}
where the coefficient is in the Laurent expansion at infinity.
The minus sign follows from $y\sim-1/\varepsilon$.
Every mixed term in the quotient is a polynomial and contributes zero.
For the scalar term, polynomial division leaves a remainder $AY+B$,
where
$A=(\phi_P(L+\ii\pi)-\phi_P(L-\ii\pi))/(2\ii\pi)$.
This proves the exact residue formula \eqref{poly:R}.

Finally, the harmonic asymptotic expansions through their $1/n$ terms
give $r_P(n)=R_P(\log n)+n^{-1}U(\log n)
+O((1+\log n)^d/n^2)$ for a polynomial $U$.
Differencing these explicit terms gives \eqref{poly:deltarasym}.
An unspecified $O(1/n)$ remainder is not being differentiated or
differenced to claim a stronger error estimate.
\end{proof}

\subsection{The complete spectral formula}

Write $\Delta p_j(n)=p_j(n)-p_j(n-1)$ and set
\begin{equation}\label{poly:JD}
 J_P=\sum_{n\ge1}\sum_{j=2}^M
          \frac{\Delta p_j(n)}{(j-1)n^{j-1}},\qquad
 D_P(s)=\sum_{n\ge1}\frac{\Delta p_1(n)}{n^s}\quad(\re s>0).
\end{equation}
The first sum is absolutely convergent. To see this directly, replace
$H_n^{(r)}$ by $H_{n-1}^{(r)}+n^{-r}$ in the finite pole polynomial;
each difference gains at least one factor $1/n$.

\begin{theorem}[Every differential-polynomial finite part]
\label{poly:bridge}
The Dirichlet series $D_P$ continues meromorphically to $\re s>-1$,
and
\begin{equation}\label{poly:bridgeeq}
 \boxed{\cI(P)=C_P+J_P+[s]D_P(s).}
\end{equation}
If $R_P'(L)=\sum a_jL^j$, an entirely ordinary-series form is
\begin{equation}\label{poly:ordinary}
 \cI(P)=C_P+J_P-\sum_j a_j\gamma_{j+1}
 -\sum_{n\ge1}
 \left\{\Delta p_1(n)-\frac{R_P'(\log n)}n\right\}\log n.
\end{equation}
The last series converges absolutely. If $\phi_P$ has degree $d\ge2$
and leading coefficient $a$, then $D_P$ has a pole of exact order
$d-1$ at zero, with leading term $-a\,d!s^{1-d}$. If
$\deg\phi_P\le1$, including $\phi_P=0$, it is holomorphic there.
\end{theorem}

\begin{proof}
The asymptotic in \eqref{poly:deltarasym} gives the normal decomposition
\[
 D_P(s)=\sum_j a_j(-1)^j\zeta^{(j)}(1+s)
 +\sum_{n\ge1}
 \left\{\Delta p_1(n)-R_P'(\log n)/n\right\}n^{-s}.
\]
Its second term is holomorphic on $\re s>-1$, including all derivatives
on compact subsets. By \eqref{poly:stieltjes}, its linear coefficient
is exactly the Stieltjes term and the last sum in
\eqref{poly:ordinary}. The leading term of $R_P'$ is
$-a d(d-1)L^{d-2}$, proving the pole classification.

It remains to establish the finite-part normalization. Integrate
\eqref{poly:MLeq}, first through a finite cutoff $N$. For $j\ge2$,
finite summation by parts gives
\[
 \sum_{n=1}^Np_j(n)(n^{1-j}-(n+1)^{1-j})
 =p_j(0)+\sum_{n=1}^N\frac{\Delta p_j(n)}{n^{j-1}}
   -\frac{p_j(N)}{(N+1)^{j-1}}.
\]
The first term cancels exactly the finite part of the pole at zero;
the last tends to zero. For simple poles the finite identity is
\[
 \sum_{n=1}^Np_1(n)\log(1+1/n)
 =p_1(N)\log(N+1)-\sum_{n=1}^N\Delta p_1(n)\log n.
\]
The boundary $p_1(N)\log(N+1)$ has logarithmic polynomial
$L R_P(L)$, with zero constant coefficient. The large-circle term
has constant $C_P$. The Stieltjes limit
$\sum_{n\le N}\log^j(n)/n
=\log^{j+1}N/(j+1)+\gamma_j+o(1)$ shows that the constant term of
the final logarithmic sum is $-[s]D_P(s)$. The subtracted series in
\eqref{poly:MLeq} is uniformly convergent after the origin pole is
removed, including on $[0,1]$, so this cutoff calculation computes the
finite part \eqref{poly:fp}. This proves both formulas.
\end{proof}

For a mixed $P$ the result simplifies to the ordinary convergent identity
\begin{equation}\label{poly:mixed}
 \boxed{\cI(P)=J_P+D_P'(0)
       =J_P-\sum_{n\ge1}\Delta p_1(n)\log n.}
\end{equation}
The disappearance of singular Laurent coordinates is thus a proved
consequence of the local cotangent algebra, not a fitted cancellation.

% END sections/02_differential_calculus.tex


% BEGIN sections/03_convergent_words.tex
\section{Convergent zeta coordinates and explicit product identities}
\label{poly:coordinates}

\subsection{A finite reduction without singular word coordinates}

Absolute convergence of a combined Dirichlet series does not by itself
make every term in an arbitrary quasi-shuffle expansion convergent.
The following construction supplies convergent coordinates individually.

\begin{theorem}[Convergent word reduction for every mixed monomial]
\label{poly:words}
Let $P$ be a mixed monomial with $K$ factors, counting multiplicities.
Then $\cI(P)$ is a finite linear combination of convergent strict
multiple-zeta values and their first spectral derivatives, of depth at
most $K-1$. The coefficients belong to the ring generated over $\Q$
by $\gamma$ and the ordinary zeta values occurring in the pole
polynomials. For $K=1$ the integral is zero. Linearity gives a finite
algorithm for every mixed polynomial.
\end{theorem}

\begin{proof}
First expand the higher-pole differences in
$h_r=H_{n-1}^{(r)}$ and $v=1/n$. A coefficient $p_j(n)$ uses at
least one singular factor, so its harmonic degree is at most $K-1$.
Taking a difference replaces at least one regular factor by a power
of $v$. Hence $\Delta p_j(n)$ has harmonic degree at most $K-2$,
and every term contains $v^q$ with $q\ge1$.
Finite stuffle multiplication replaces products of the $h_r$ by strict
words. Therefore
\[
 \sum_{n\ge1}\frac{v^q H_{n-1}(w)}{n^{j-1}}
 =Z(q+j-1,w),\qquad j\ge2,
\]
with first entry at least two and depth at most $K-1$.
This gives the claimed representation of $J_P$.

For the residue, express $p_1(n)$ as a polynomial
$F(B,A_2,A_3,\ldots)$ with $B=H_n-\gamma$ and
$A_r=A_r(n)$. In \eqref{poly:localpsi} this amounts to replacing
$c_{r,n}$ by $(1+(-1)^r)\zeta(r)-A_r$.
The residue cancellation theorem, with $B$ an indeterminate, proves
the polynomial identity $F(B,0,0,\ldots)=0$. Thus every monomial
of $F$ contains at least one upper tail.

Put $b=H_{n-1}-\gamma$ and use the exact increment formula
\begin{equation}\label{poly:taildifference}
 \Delta p_1(n)=F(b+v,A_2,A_3,\ldots)
       -F(b,A_2+v^2,A_3+v^3,\ldots).
\end{equation}
Each resulting term is a constant times
\begin{equation}\label{poly:splitword}
 n^{-q}H_{n-1}(w)\,U_n(v_1,\ldots,v_d),\qquad
 U_n(v_1,\ldots,v_d)=
 \sum_{k_1>\cdots>k_d>n}\prod k_j^{-v_j}.
\end{equation}
Here all upper letters $v_j$ are at least two. Stuffle multiplication
has been applied separately above and below the distinguished index $n$.
If no upper letter remains, at least one removed tail supplies
$n^{-r}$ with $r\ge2$, so $q\ge2$. If an upper letter remains,
$q\ge1$ is sufficient. The corresponding sums and logarithmic sums are
\begin{align}
 \sum_{n\ge1}n^{-q}H_{n-1}(w)U_n(v)
     &=Z(v_1,\ldots,v_d,q,w),\label{poly:concatenate}\\
 -\sum_{n\ge1}n^{-q}H_{n-1}(w)U_n(v)\log n
     &=\partial_{d+1}Z(v_1,\ldots,v_d,q,w).
 \label{poly:concatenateD}
\end{align}
Every displayed word has first entry at least two. Its derivative is
absolutely convergent as well: a logarithm is dominated by an arbitrarily
small extra power of the largest index, while the standard admissible
series has a positive convergence margin.
The difference in \eqref{poly:taildifference} removes at least one of
at most $K-1$ harmonic or tail factors; reinserting the distinguished
index gives depth at most $K-1$. Formula \eqref{poly:mixed} completes
the proof. For one positive-order polygamma all its principal coefficients
are constant in $n$ and its simple residue is zero, giving $J_P=D_P=0$.
\end{proof}

This theorem makes a finite-coordinate statement. It does not assert
that its resulting multiple-zeta derivatives reduce to ordinary zeta
derivatives, nor that distinct remaining coordinates are independent.
For polynomial coefficients already containing other constants, those
constants are simply retained as coefficients.

\subsection{Polylogarithmic kernels for the same coordinates}

Let
\begin{equation}\label{poly:S}
 S_j(z)=\sum_{n\ge1}\Delta p_j(n)z^n,\qquad |z|<1.
\end{equation}
The finite harmonic expansion shows that every $S_j$ is a finite linear
combination of strict multiple polylogarithms:
\[
 \sum_{n\ge1}\frac{H_{n-1}(w)}{n^q}z^n=\Li_{q,w}(z),
 \qquad
 \Li_{s_1,\ldots,s_d}(z)=
 \sum_{n_1>\cdots>n_d\ge1}\frac{z^{n_1}}{n_1^{s_1}\cdots n_d^{s_d}}.
\]
For a mixed polynomial, $S_1$ has the continuous endpoint value
$S_1(1)=\sum\Delta p_1(n)=-p_1(0)$.

\begin{corollary}[Ordinary logarithmic polylogarithm bridge]
\label{poly:polylog}
For mixed $P$,
\begin{equation}\label{poly:polylogeq}
 \boxed{\displaystyle
 \cI(P)=J_P+\int_0^1
 \frac{S_1(z)-zS_1(1)}{z(-\log z)}\dd z.}
\end{equation}
The integral is ordinary and absolutely convergent. Moreover,
\begin{equation}\label{poly:Jpolylog}
 J_P=\sum_{j=2}^M\frac1{(j-1)!}
 \int_0^1\frac{(-\log z)^{j-2}}z S_j(z)\dd z.
\end{equation}
\end{corollary}
\begin{proof}
The estimate $\Delta p_1(n)=O(n^{-2}(1+\log n)^d)$ implies
$\sum |\Delta p_1(n)|\log n<\infty$. Frullani's identity in the
form
\[
 \int_0^1\frac{z^{n-1}-1}{-\log z}\dd z=-\log n
\]
therefore justifies termwise integration of the complete subtracted
kernel, including absolute integrability. This gives $D_P'(0)$.
For the second formula use
$\int_0^1z^{n-1}(-\log z)^{j-2}\dd z=(j-2)!/n^{j-1}$
and the absolute convergence of the harmonic differences.
Individual polylogarithms in a representation of $S_1$ can diverge at
one; $S_1(1)$ means the value of their complete combination, not a
termwise substitution into a divergent expression.
\end{proof}

\subsection{Every bilinear polygamma integral}

The entire bilinear layer collapses to two ordinary zeta coordinates.
This is also a useful explicit check on the signs and endpoint
normalization in the general theorem.

\begin{theorem}[Bilinear evaluation]
\label{poly:bilinear}
For nonnegative integers $r,s$ with $r+s\ge1$, put $h=r+s$,
$\alpha=(-1)^{r+1}$, and $\beta=(-1)^{s+1}$. Then
\begin{align}\label{poly:bilineareq}
 \boxed{\displaystyle
 \FP\int_0^1\psi^{(r)}(x)\psi^{(s)}(x)\dd x
 =h!\left\{(\alpha+\beta)\zeta'(h+1)
 +\bigl[\alpha(H_h-H_s)+\beta(H_h-H_r)\bigr]\zeta(h+1)\right\}.}
\end{align}
In particular, for $r\ge1$,
\begin{equation}\label{poly:square}
 \FP\int_0^1\bigl(\psi^{(r)}(x)\bigr)^2\dd x
 =2(-1)^{r+1}(2r)!
 \{\zeta'(2r+1)+(H_{2r}-H_r)\zeta(2r+1)\}.
\end{equation}
\end{theorem}

\begin{proof}
Assume first $r,s\ge1$. Multiplication of the two differentiated
series in \eqref{poly:localpsi} gives
\[
 p_1(n)=(\alpha+\beta)h!c_{h+1,n}.
\]
The parity of $h$ shows that either $\alpha+\beta=0$ or
$h+1$ is odd; in the latter case $c_{h+1,n}=-A_{h+1}(n)$.
Thus
$D_P(s_0)=(\alpha+\beta)h!\zeta(h+1+s_0)$.
The higher-pole differences give
\[
 J_P=\left\{
 \alpha r!\sum_{k=1}^r\frac{(h-k)!}{(r-k)!k}
 +\beta s!\sum_{k=1}^s\frac{(h-k)!}{(s-k)!k}
 \right\}\zeta(h+1).
\]
The finite binomial identity
\begin{equation}\label{poly:finiteharmonic}
 r!\sum_{k=1}^r\frac{(r+s-k)!}{(r-k)!k}
 =(r+s)!(H_{r+s}-H_s)
\end{equation}
follows, for example, by taking the logarithmic derivative in $a$ of
Vandermonde's identity for $[x^r](1-x)^{-a}(1-x)^{-s-1}$ at
$a=0$, or by induction in $r$. To see the coefficient calculation
directly, $\left.\partial_a(a)_k/k!\right|_{a=0}=1/k$ for $k\ge1$,
and the derivative of $(a+s+1)_r/r!$ is
$\binom{r+s}{r}(H_{r+s}-H_s)$. Multiplication by $r!s!$
gives \eqref{poly:finiteharmonic}.

If $r=0<s$, the coefficient of the pole of order $s+1$ contains
$\beta s!b_n$ instead. Its difference is $\beta s!/n$, and it
supplies exactly the last $k=s$ term of the same finite sum.
The simple-pole formula remains valid. The case $s=0<r$ is symmetric.
Formula \eqref{poly:mixed} now gives the result.
\end{proof}

For example,
\begin{align}
 \FP\int_0^1(\psi')^2\dd x&=2\zeta(3)+4\zeta'(3),\label{poly:trigamma2}\\
 \FP\int_0^1\psi\psi''\dd x&=-3\zeta(3)-4\zeta'(3),\\
 \FP\int_0^1\psi'\psi''\dd x&=-3\zeta(4),\\
 \FP\int_0^1(\psi'')^2\dd x&=-28\zeta(5)-48\zeta'(5),\\
 \FP\int_0^1\psi'\psi'''\dd x&=32\zeta(5)+48\zeta'(5).
\end{align}
If $r$ and $s$ have opposite parity, the zeta derivative disappears
exactly. The exceptional pair $(0,0)$ has a singular scalar residue
and instead gives $\cI(u_0^2)=2\gamma_1-2\zeta(2)$ by
Theorem~\ref{poly:bridge}.

% END sections/03_convergent_words.tex


% BEGIN sections/04_normal_forms.tex
\section{Canonical integration by parts with endpoint constants}
\label{poly:normalform}

Let $\mathcal A=\Q[u_0,u_1,u_2,\ldots]$ and
\[
 \cD=\sum_{r\ge0}u_{r+1}\frac{\partial}{\partial u_r},\qquad
 \wt(u_r)=r+1.
\]
Each polynomial uses finitely many variables. The weight-$w$ subspace
$\mathcal A_w$ has dimension $p(w)$, the partition number.
Let $\mathcal B_w$ consist of $u_0^w$ and all weight-$w$ monomials
whose highest positive-index variable occurs at least twice.
For $w=0$ use $\mathcal B_0=\{1\}$.

\begin{theorem}[A triangular normal form at every weight]
\label{poly:IBP}
For every $w\ge2$ there is a direct sum
\begin{equation}\label{poly:directsum}
 \mathcal A_w=\cD\mathcal A_{w-1}\oplus
                      \operatorname{span}_{\Q}\mathcal B_w.
\end{equation}
Every $P\in\mathcal A_w$ therefore has a unique decomposition
$P=\cD Q+N$ with $Q\in\mathcal A_{w-1}$ and $N$ in the displayed normal space. In particular
\begin{equation}\label{poly:dimension}
 |\mathcal B_w|=p(w)-p(w-1)\quad(w\ge2).
\end{equation}
The quotient has dimension one at weights zero and one, so its
Hilbert series is
\begin{equation}\label{poly:Hilbert}
 t+\prod_{j\ge2}(1-t^j)^{-1}.
\end{equation}
\end{theorem}

\begin{proof}
Suppose a monomial has highest variable $u_r$, $r\ge1$, occurring
exactly once. Write it as $u_r u_{r-1}^a V$, where $V$ involves
only $u_0,\ldots,u_{r-2}$. Then
\begin{equation}\label{poly:IBPstep}
 u_r u_{r-1}^aV=
 \cD\left(\frac{u_{r-1}^{a+1}V}{a+1}\right)
 -\frac{u_{r-1}^{a+1}\cD V}{a+1}.
\end{equation}
Every monomial in the last term has strictly smaller highest index.
Iterating therefore terminates and proves that the two spaces span.

For directness, let $Q$ be a nonconstant polynomial with highest
variable $u_m$. The highest-index terms of $\cD Q$ are
$u_{m+1}\partial Q/\partial u_m$, all with their highest variable
occurring once. They cannot occur in a polynomial in the normal space.
Since the characteristic is zero, $\partial Q/\partial u_m\ne0$.
This also proves that $\cD$ is injective on every positive-weight
subspace. Taking dimensions proves \eqref{poly:dimension}.
The weight-zero kernel consists of constants, which explains the extra
$t$ in \eqref{poly:Hilbert}.
\end{proof}

\begin{center}
\small
\begin{tabular}{ccl}
\toprule
Weight & Dimension & A canonical basis, with $u_r=\psi^{(r)}$\\
\midrule
2&1&$u_0^2$\\
3&1&$u_0^3$\\
4&2&$u_0^4,\ u_1^2$\\
5&2&$u_0^5,\ u_0u_1^2$\\
6&4&$u_0^6,\ u_0^2u_1^2,\ u_1^3,\ u_2^2$\\
7&4&$u_0^7,\ u_0^3u_1^2,\ u_0u_1^3,\ u_0u_2^2$\\
8&7&$u_0^8,\ u_0^4u_1^2,\ u_0^2u_1^3,\ u_1^4,
 u_0^2u_2^2,\ u_1u_2^2,\ u_3^2$\\
\bottomrule
\end{tabular}
\end{center}

\subsection{The endpoint map is part of the theorem}

For a polynomial $Q$, define
\begin{equation}\label{poly:endpoint}
 E(Q)=f_Q(1)-[x^0]f_Q(x).
\end{equation}
Both terms are finite polynomial calculations. The values at one are
$\psi(1)=-\gamma$ and
$\psi^{(r)}(1)=(-1)^{r+1}r!\zeta(r+1)$ for $r\ge1$.
The constant at zero is obtained from \eqref{poly:localpsi} with $n=0$;
it is not obtained merely by multiplying the individual constant terms.

\begin{corollary}[Exact finite-part integration by parts]
\label{poly:endpointcor}
For the normal decomposition $P=\cD Q+N$,
\begin{equation}\label{poly:IBPintegral}
 \boxed{\cI(P)=E(Q)+\cI(N).}
\end{equation}
For a common shift $a$ with $\re a>0$, the corresponding ordinary
identity is
\begin{equation}\label{poly:IBPshift}
 \int_0^1f_P(a+x)\dd x
 =f_Q(a+1)-f_Q(a)+\int_0^1f_N(a+x)\dd x.
\end{equation}
\end{corollary}
\begin{proof}
On $[\varepsilon,1]$ the integral of $f_{\cD Q}$ is
$f_Q(1)-f_Q(\varepsilon)$. Extracting its original-coordinate
finite part retains precisely $f_Q(1)-[x^0]f_Q(x)$.
The shifted identity is the ordinary fundamental theorem of calculus.
\end{proof}

For example, $\cD(u_0u_1)=u_1^2+u_0u_2$, but
\begin{equation}\label{poly:boundaryexample}
 E(u_0u_1)=-\zeta(3),\qquad
 \cI(u_1^2)+\cI(u_0u_2)=-\zeta(3).
\end{equation}
Likewise $E(u_1u_2)=4\zeta(5)$. These agree with all five bilinear
specializations above. Dropping total derivatives would lose these
constants and give a false relation.

The dimension theorem answers the stated formal classification problem
before imposing arithmetic identities among zeta coordinates.
Its numerical image may have smaller dimension. For instance the
$u_2^2$ coordinate at weight six already has the elementary zeta-derivative
evaluation \eqref{poly:square}. Neither the formal normal form nor a
failure to find further numerical relations proves arithmetic independence.

% END sections/04_normal_forms.tex


% BEGIN sections/05_primitives.tex
\section{Normalized primitives and common shifted endpoints}
\label{poly:primitives}

The global expansion also answers the endpoint-normalization part of
the common-shift question. It yields convergent formulas at independent
arguments without replacing ordinary harmonic numbers by shifted ones
without justification.

For $j\ge1$ put
\[
 F_1(z)=\Log z,\qquad F_j(z)=\frac{z^{1-j}}{1-j}\quad(j\ge2).
\]
With the data of \eqref{poly:poles}, define
\begin{align}\label{poly:primitive}
 \mathcal P_P(z)={}&\sum_{j=1}^Mp_j(0)(F_j(z)-F_j(1))
                  +\kappa_P(z-1)\notag\\
 &+\sum_{n\ge1}\sum_{j=1}^Mp_j(n)
 \left\{F_j(n+z)-F_j(n+1)-\frac{z-1}{n^j}\right\}.
\end{align}

\begin{theorem}[Primitive and shift theorem]
\label{poly:primitivethm}
The function $\mathcal P_P$ is holomorphic on
$\C\setminus(-\infty,0]$, its series converges normally there,
and
\begin{equation}\label{poly:primitiveidentity}
 \mathcal P_P'(z)=f_P(z),\qquad \mathcal P_P(1)=0.
\end{equation}
In particular, for $\re a>0$,
\begin{equation}\label{poly:shift}
 I_P(a):=\int_0^1f_P(a+x)\dd x
        =\mathcal P_P(a+1)-\mathcal P_P(a).
\end{equation}
As $a\to0$ in any fixed closed sector $|\arg a|\le\pi-\delta$,
\begin{equation}\label{poly:shiftlimit}
 I_P(a)=\sum_{j=2}^M\frac{p_j(0)}{j-1}a^{1-j}
          -p_1(0)\Log a+\cI(P)+O(a).
\end{equation}
The last expansion refers to analytic continuation of the difference
in \eqref{poly:shift}; on the positive real axis it is an ordinary
integral. Thus subtracting its displayed singular terms recovers
exactly the finite part at zero.
\end{theorem}

\begin{proof}
For $z$ in a compact set, the braces in \eqref{poly:primitive}
are $O(n^{-j-1})$ by integrating
$(n+z)^{-j}-n^{-j}$ from $1$ to $z$. The pole coefficients have only
polylogarithmic growth, so the series and its derivative converge
normally. The derivative is \eqref{poly:MLeq}, and every term vanishes
at one. This fixes the integration constant.
Integrating along the horizontal segment from $a$ to $a+1$ proves
\eqref{poly:shift}. Near zero, integrate the convergent Laurent expansion
of $f_P$; its analytic remainder has an $O(a)$ primitive. The upper
endpoint contributes another $O(a)$ term. The remaining constant is
the definition \eqref{poly:fp}. The local convergent expansion also
proves uniformity on the stated sectors.
\end{proof}

There is an exact differential check using only the Gamma recurrence:
\begin{align}\label{poly:shiftdifferential}
 I_P'(a)&=P\left(\psi(a)+a^{-1},
 \psi'(a)-a^{-2},\ldots,
 \psi^{(r)}(a)+(-1)^rr!a^{-r-1},\ldots\right)
               -f_P(a).
\end{align}
This follows from $f_P(a+1)-f_P(a)$ and is valid throughout
$\re a>0$. The explicit endpoint law is the part settled here.
A further reduction of all such shifted logarithmic sums to a specified
common-shift multiple-Hurwitz basis is not claimed.

\subsection{A multiple-polylogarithm representation of the primitives}

Put $V_j(z)=\sum_{n\ge1}p_j(n)z^n$ for $|z|<1$. Finite stuffle
expansion of the inclusive harmonic polynomials expresses $V_j$ using
multiple polylogarithms and rational factors $1/(1-z)$.
For example, for a nonempty word $w$,
$\sum_{n\ge1}H_n(w)z^n=\Li_w(z)/(1-z)$; constants have
generating function $z/(1-z)$.

For $\re z>0$, the expansion and primitive have ordinary kernels
\begin{align}\label{poly:MLkernel}
 f_P(z)={}&P_0(z)+\kappa_P+
 \sum_{j=1}^M\frac1{(j-1)!}
 \int_0^\infty t^{j-1}(e^{-zt}-1)V_j(e^{-t})\dd t,\\
 \mathcal P_P(z)={}&\sum_{j=1}^Mp_j(0)(F_j(z)-F_j(1))
 +\kappa_P(z-1)\notag\\
 &+\sum_{j=1}^M\frac1{(j-1)!}
 \int_0^\infty t^{j-1}
 \left\{\frac{e^{-t}-e^{-zt}}t-(z-1)\right\}
 V_j(e^{-t})\dd t.\label{poly:primitivekernel}
\end{align}
Near $t=0$, $V_j(e^{-t})=O(t^{-1}(1+|\log t|)^d)$, while the
braced differences vanish to first order in $t$. Near infinity the
positive summation indices supply exponential decay. These estimates,
uniform on compact subsets of $\re z>0$, justify Fubini and all fixed
$z$-derivatives. The normal partial-fraction formula gives the continuation
to the larger slit plane without imposing an unjustified integral
convergence claim there.

This construction links three operations under one normalization:
spatial differentiation multiplies the Mellin kernel by powers of $-t$;
spectral differentiation of the harmonic Dirichlet series inserts
logarithms of the distinguished integer; and spatial integration gives
the explicit quotient in \eqref{poly:primitivekernel}. These operations
are distinct and their constants have been fixed separately.

% END sections/05_primitives.tex


% BEGIN sections/06_appell.tex
\section{One singular Laurent coordinate at every Gamma--Appell degree}
\label{ac:section}

The newest canonical-digamma report proves that a Gamma--Appell sequence
matches the nondecaying harmonic residues of the digamma powers, and
asks which lower spectral coordinates survive
\cite{Fifth}. The higher-power and mixed-resolvent reports give
terminating harmonic-word compilers and complete fifth-power formulas,
but leave the sixth and seventh cases and a uniform preferred spectral
support as explicit further targets \cite{Higher,Mixed}.
The result here settles a precise part of that question at every degree:
one normalization removes \emph{all lower singular Laurent coordinates}.
The remaining terms have ordinary absolutely convergent multiple-zeta
derivatives. It does not remove every lower-depth arithmetic constant,
and it makes no assertion of arithmetic independence.

\subsection{The normalized moments and their finite data}

Define
\begin{equation}\label{ac:appell}
 \mathsf P_m(Y)=m![t^m]\Gamma(1-t)e^{tY},\qquad
 \mathsf S_m=\FP\int_0^1\mathsf P_m(\psi(x))\,\dd x.
\end{equation}
The finite part uses the original cutoff $x=\varepsilon$.
For $m\ge2$, write
\begin{equation}\label{ac:coordinates}
 \frac1{\Gamma(1+t)}=\sum_{j\ge0}c_jt^j,
 \qquad
 \lambda_d=[s^1]Z_d(1+s,\underbrace{1,\ldots,1}_{d-1}),
 \quad d\ge1.
\end{equation}
Here $Z_d$ is the strict multiple zeta function, with the largest
summation index first. The coefficient of a Laurent germ has no
additional factorial or alternating sign. In particular
$\lambda_1=-\gamma_1$.

Let $p_{m,j}(n)$ be the coefficient of $(x+n)^{-j}$ in
$\mathsf P_m(\psi(x))$ at $x=-n$, and put
\begin{align}
 e_r(n)&=H_n(\{1\}^r),\qquad e_0(n)=1,\notag\\
 \delta_m(n)&=p_{m,1}(n)+m!e_{m-1}(n),\label{ac:defect}\\
 \mathsf J_m&=\sum_{n\ge1}\sum_{j=2}^m
       \frac{p_{m,j}(n)-p_{m,j}(n-1)}{(j-1)n^{j-1}},
       \label{ac:J}\\
 \mathsf D_m&=-\sum_{n\ge1}
                \bigl(\delta_m(n)-\delta_m(n-1)\bigr)\log n.
       \label{ac:D}
\end{align}
The infinite expressions in \eqref{ac:J} and \eqref{ac:D} will be
replaced by finite multiple-zeta formulas. They are also useful
ordinary convergent representatives for numerical evaluation.

\begin{theorem}[All lower singular Laurent coordinates disappear]
\label{ac:main}
For every integer $m\ge2$,
\begin{equation}\label{ac:mainformula}
 \boxed{\displaystyle
 \mathsf S_m=m!\bigl(c_m-\lambda_{m-1}\bigr)
                 +\mathsf J_m+\mathsf D_m.}
\end{equation}
The following assertions specify a finite arithmetic normal form.
\begin{enumerate}
 \item $\mathsf J_m$ is a finite linear combination of ordinary
 admissible multiple zeta values, with coefficients in
 $\Q[\zeta(2),\ldots,\zeta(m)]$.
 \item $\mathsf D_2=\mathsf D_3=0$. For $m\ge4$, $\mathsf D_m$
 is a finite linear combination of first partial derivatives
 $\partial_j Z_d(s_1,\ldots,s_d)$ at admissible positive-integer
 indices, with $d\le m-3$. Each derivative has an ordinary
 absolutely convergent defining series in a full parameter
 neighborhood of its evaluation point.
 \item Every fixed word in these finite expressions has weight
 at most $m-1$. Assigning weight $r$ to a coefficient $\zeta(r)$
 makes each summand have total weight $m-1$ before the spectral
 differentiation.
\end{enumerate}
Thus $\lambda_{m-1}$ is the only singular Laurent coefficient
needed by this construction. No lower $\lambda_d$ is retained.
\end{theorem}

\begin{proof}[The finite-part identity]
The scalar version of Theorem~\ref{poly:bridge} is linear in its
polynomial argument. Its large-circle constant for $\mathsf P_m$ is
\[
 m![t^m]\Gamma(1-t)\frac1{2\pi}
       \int_{-\pi}^{\pi}e^{it\theta}\,\dd\theta
 =m![t^m]\Gamma(1-t)\frac{\sin\pi t}{\pi t}
 =m!c_m.
\]
The last equality is Euler's reflection formula. The generating
identity
\[
 \operatorname*{Res}_{x=-n}
 \frac{\Gamma(1-t)\Gamma(x)}{\Gamma(x-t)}
       =-t\prod_{k=1}^n(1+t/k)
\]
explains the comparison term $-m!e_{m-1}(n)$; its use here is
consistent with the established differential normalization in
\cite{Fifth}. Since
\[
 e_{m-1}(n)-e_{m-1}(n-1)=\frac{e_{m-2}(n-1)}n,
\]
the meromorphic residue-difference series is
\begin{equation}\label{ac:Dsplit}
 \sum_{n\ge1}\frac{p_{m,1}(n)-p_{m,1}(n-1)}{n^s}
 =-m!Z_{m-1}(1+s,\{1\}^{m-2})
   +\sum_{n\ge1}\frac{\delta_m(n)-\delta_m(n-1)}{n^s}.
\end{equation}
The next proof shows that the second series is holomorphic for
$\Re s>-1$ and that its coefficient of $s$ is exactly
$\mathsf D_m$. Extracting the coefficient in the scalar bridge
gives \eqref{ac:mainformula}, including its endpoint and angular
constants. The finite-word description of $\mathsf J_m$ follows
from the same finite-difference and quasi-shuffle construction as
Theorem~\ref{poly:words}; its initial index is at least two.
\end{proof}

\subsection{Why the correction is a finite convergent derivative family}

The proof can be made completely explicit without assigning a new
regularization to any mixed term. Set
\begin{equation}\label{ac:tails}
 h=H_n,\qquad A_r(n)=\zeta(r)-H_n^{(r)}=\sum_{k>n}k^{-r},
                     \qquad r\ge2.
\end{equation}

\begin{lemma}[The tail ideal and the depth bound]
\label{ac:tailideal}
For every fixed $m$, $\delta_m(n)$ is a polynomial
$\delta_m(h,A_2,A_3,\ldots)$ with coefficients in the ordinary
zeta algebra. It vanishes identically when all $A_r$ are set to
zero. For $m\ge4$, its degree in the harmonic variables
$h,A_2,A_3,\ldots$, counting each of them with degree one, is at
most $m-3$. Moreover $\delta_2=\delta_3=0$.
\end{lemma}

\begin{proof}
The exact local Laurent expansion is
\begin{equation}\label{ac:local}
 \psi(-n+v)=H_n-\gamma-\frac1v
       +\sum_{r\ge1}\bigl(H_n^{(r+1)}+(-1)^{r+1}\zeta(r+1)\bigr)v^r.
\end{equation}
Only finitely many coefficients enter any fixed residue. Substituting
$H_n^{(r)}=\zeta(r)-A_r$ proves polynomiality.
With every $A_r=0$, the right side is
$h-\gamma-\pi\cot\pi v$. The universal cotangent residue
\[
 \Res_{v=0}e^{-\pi t\cot\pi v}=-\frac{\sin\pi t}{\pi}
\]
therefore gives the zero-tail residue generator (the residue identity follows coefficientwise from \eqref{poly:reschange})
\[
 -\Gamma(1-t)e^{(h-\gamma)t}\frac{\sin\pi t}{\pi}
 =-t\frac{e^{(h-\gamma)t}}{\Gamma(1+t)}.
\]
On the other hand the zero-tail specialization of the elementary
harmonic generator is
\[
 \sum_{r\ge0}e_r t^r
 =\exp\left(ht+\sum_{r\ge2}\frac{(-1)^{r-1}\zeta(r)}r t^r\right)
 =\frac{e^{(h-\gamma)t}}{\Gamma(1+t)}.
\]
Their sum in \eqref{ac:defect} is zero. This proves membership in
the ideal generated by the $A_r$, using exact Gamma identities.

For the harmonic degree, first keep $h_r=H_n^{(r)}$ as independent
variables and write $h=h_1$. The centered Gamma--Appell generator is
\[
 \Gamma(1-t)e^{tY}
 =\exp\left((Y+\gamma)t+\sum_{r\ge2}\frac{\zeta(r)}r t^r\right).
\]
For $m\ge3$, the only terms of harmonic degree $m-1$ and $m-2$
in its simple-pole coefficient are respectively
\[
 -m h^{m-1},\qquad
 \frac{m!}{2(m-3)!}h_2h^{m-3}.
\]
Indeed the first chooses one $-1/v$. To have only one fewer
harmonic factor, one must choose two $-1/v$ factors and one
$h_2v$ factor. Any other positive power of $v$, additional
negative-power factor, or ordinary-zeta coefficient lowers the
harmonic degree further. The same two terms occur in
$-m!e_{m-1}$ by expansion of the elementary harmonic generator.
They cancel in $\delta_m$. Replacing $h_r$ by $\zeta(r)-A_r$
cannot increase this degree. Direct extraction at $m=2,3$ gives
zero defects, completing the proof.
\end{proof}

Let $h=H_{n-1}$ and $v=1/n$, while $A_r=A_r(n)$ is the
\emph{strict upper} tail. The finite difference is consequently
the exact polynomial
\begin{equation}\label{ac:delta-rule}
 \delta_m(n)-\delta_m(n-1)
 =\delta_m(h+v,A_2,A_3,\ldots)
   -\delta_m(h,A_2+v^2,A_3+v^3,\ldots).
\end{equation}
Every monomial has the form
\begin{equation}\label{ac:monomial}
 a\,v^q h^b\prod_{r\ge2}A_r^{b_r},\qquad q\ge1.
\end{equation}
If it contains no upper tail, then $q\ge2$. This last restriction
is decisive: a term with $q=1$ arises only by changing $h$, so
the tail-ideal property leaves an upper tail behind.

For a word $\mathbf a=(a_1,\ldots,a_r)$ of entries at least two,
define
\[
 T_{\mathbf a}(n)=\sum_{k_1>\cdots>k_r>n}
                        k_1^{-a_1}\cdots k_r^{-a_r}.
\]
The empty upper word and empty lower word both have value one.
Apply the strict quasi-shuffle rule separately to the upper-tail
product and to $h^b$. Each monomial becomes a finite sum of
$n^{-q}T_{\mathbf a}(n)H_{n-1}(\mathbf b)$. Then
\begin{equation}\label{ac:concatenate}
 \sum_{n\ge1}\frac{T_{\mathbf a}(n)H_{n-1}(\mathbf b)}{n^{q+s}}
       =Z(\mathbf a,q+s,\mathbf b).
\end{equation}
The index ordering on the two sides is identical. If the upper word
is empty then $q\ge2$; otherwise the first entry is at least two.
All entries are positive integers at $s=0$, so these series and
their first derivatives converge absolutely on a small complex
neighborhood. Equivalently, elementary tail estimates in
\eqref{ac:monomial} give $O(n^{-2}(1+\log n)^B)$ before
inserting $n^{-s}$, proving normal convergence for $\Re s>-1$.

Differentiation of \eqref{ac:concatenate} therefore gives the
explicit conversion rule
\begin{equation}\label{ac:marked-rule}
 -\sum_{n\ge1}\frac{T_{\mathbf a}(n)H_{n-1}(\mathbf b)\log n}{n^q}
 =\partial_{r+1}Z(\mathbf a,q,\mathbf b),\qquad r=|\mathbf a|.
\end{equation}
Taking a nonzero finite difference replaces at least one harmonic
factor by a power of $v$. The remaining harmonic-factor count is
at most $m-4$; adjoining the distinguished $n$ in
\eqref{ac:concatenate} gives depth at most $m-3$.
This completes the proof of Theorem~\ref{ac:main} and provides a
terminating finite algorithm for every term it retains.

\subsection{Explicit sixth and seventh normalized powers}

For an admissible integer word $\mathbf w$, write
$Z_{\mathbf w}=Z(\mathbf w)$ and let
\begin{equation}\label{ac:directions}
 \Delta_{ij}Z_{\mathbf w}=(\partial_i-\partial_j)Z(\mathbf w).
\end{equation}
These are derivatives of a genuinely free-order analytic function.
They are not obtained by differentiating an identity valid only
at integer orders.

\begin{theorem}[The sixth normalized digamma moment]
\label{ac:sixth}
The sixth-degree identity is
\begin{equation}\label{ac:sixthmain}
 \boxed{\displaystyle
 \mathsf S_6=720(c_6-\lambda_5)
                +245\zeta(2)\zeta(3)+175\zeta(5)+\mathsf D_6,}
\end{equation}
where the entire mixed correction is the following finite sum of
convergent derivatives:
\begin{align}\label{ac:D6}
 \mathsf D_6={}&120\zeta(2)\Delta_{12}Z_{2,1}
      -60\Delta_{23}Z_{2,2,1}
      +120\Delta_{12}Z_{3,1,1}\notag\\
 &-180\Delta_{12}Z_{4,1}-30\Delta_{12}Z_{2,3}
      +138\zeta'(5)-60\zeta(2)\zeta'(3)
                        +30\zeta(3)\zeta'(2).
\end{align}
\end{theorem}

\begin{theorem}[The seventh normalized digamma moment]
\label{ac:seventh}
The seventh-degree identity is
\begin{equation}\label{ac:seventhmain}
 \boxed{\displaystyle
 \mathsf S_7=5040(c_7-\lambda_6)
       +\frac{3731}{4}\zeta(3)^2
       +\frac{918631}{240}\zeta(6)+\mathsf D_7,}
\end{equation}
where
\begin{align}\label{ac:D7}
 \mathsf D_7={}&840\zeta(2)\Delta_{12}Z_{2,1,1}
       -840\zeta(2)\Delta_{12}Z_{3,1}
       +210\zeta(2)\partial_1 Z_{2,2}\notag\\
 &-420\Delta_{23}Z_{2,2,1,1}
       +630\Delta_{23}Z_{2,3,1}
       +840\Delta_{12}Z_{3,1,1,1}\notag\\
 &+420\Delta_{12}Z_{3,1,2}
       +210(2\Delta_{13}+\Delta_{23})Z_{3,2,1}
       +210\partial_1 Z_{2,2,2}\notag\\
 &+798\partial_1 Z_{3,3}
       -1260\Delta_{12}Z_{4,1,1}
       +168\Delta_{12}Z_{2,4}+1596\Delta_{12}Z_{5,1}\notag\\
 &+\frac{651}{2}\zeta(4)\zeta'(2)
       +336\zeta(2)\zeta'(4)
       -378\zeta(3)\zeta'(3)-504\zeta'(6).
\end{align}
\end{theorem}

For completely explicit integrands, put $U=\psi(x)+\gamma$.
Then
\begin{align}\label{ac:P6}
 \mathsf P_6(\psi)={}&U^6+15\zeta(2)U^4+40\zeta(3)U^3
       +\frac{405}{2}\zeta(4)U^2\notag\\
 &+(120\zeta(2)\zeta(3)+144\zeta(5))U
       +40\zeta(3)^2+\frac{2745}{8}\zeta(6),\\
 \mathsf P_7(\psi)={}&U^7+21\zeta(2)U^5+70\zeta(3)U^4
       +\frac{945}{2}\zeta(4)U^3\notag\\
 &+(420\zeta(2)\zeta(3)+504\zeta(5))U^2
       +\left(280\zeta(3)^2+\frac{19215}{8}\zeta(6)\right)U\notag\\
 &+945\zeta(3)\zeta(4)+504\zeta(2)\zeta(5)+720\zeta(7).
       \label{ac:P7}
\end{align}
% Suggested replacement/extension of the paragraph following ac:P7.
% Uses only the macros and citation keys already used by appell_section.tex.
\subsection{Recovering raw moments without external tables}

Put $Q_m=\FP\int_0^1\psi(x)^m\,\dd x$, with $Q_0=1$.
The inverse of the Gamma--Appell normalization is the finite formula
\begin{equation}\label{ac:rawclosed}
 \boxed{\displaystyle
 Q_m=m!\sum_{k=0}^m\frac{(-1)^{m-k}c_{m-k}}{k!}\,\mathsf S_k.}
\end{equation}
Indeed, the defining polynomial identity gives the formal-power-series
relation
\[
 \sum_{m\ge0}\frac{\mathsf S_m}{m!}t^m
   =\Gamma(1-t)\sum_{m\ge0}\frac{Q_m}{m!}t^m.
\]
Multiplication by $1/\Gamma(1-t)=\sum_{j\ge0}(-1)^jc_jt^j$
proves \eqref{ac:rawclosed}. This is a coefficientwise polynomial
argument; it requires no interchange of an infinite series with a
finite-part integral. The coefficients themselves are determined by
\begin{equation}\label{ac:crecurrence}
 c_0=1,\qquad
 j c_j=\gamma c_{j-1}
       +\sum_{r=2}^j(-1)^{r+1}\zeta(r)c_{j-r},\qquad j\ge1.
\end{equation}

For completeness, the lower degrees already treated in the earlier
reports \cite{Higher,Fifth} take the following form in the
present normalization. They are recovered prior inputs, included to
make the sixth- and seventh-power evaluations self-contained:
\begin{align}
 \mathsf S_0&=1,\qquad \mathsf S_1=\gamma,
       \label{ac:lowerinitial}\\
 \mathsf S_2&=2(c_2-\lambda_1),\nonumber\\
 \mathsf S_3&=6(c_3-\lambda_2)+3\zeta(2),\nonumber\\
 \mathsf S_4&=24(c_4-\lambda_3)+12\zeta(3)+4\zeta'(3),\nonumber\\
 \mathsf S_5&=120(c_5-\lambda_4)+\frac{260}{3}\zeta(4)
                     +\mathsf D_5,\label{ac:lowernormalized}\\
 \mathsf D_5&=20\zeta(2)\zeta'(2)-10\partial_2Z_{2,2}
       +20\Delta_{12}Z_{3,1}-30\zeta'(4).
       \label{ac:lowerD5}
\end{align}
Here $Z_{a,b}=Z_2(a,b)$ and
$\Delta_{12}=\partial_1-\partial_2$, with all derivatives evaluated
at the displayed indices. In particular, all terms in
$\mathsf D_5$ have ordinary absolutely convergent defining series.
The value $\mathsf S_1=\gamma$ uses
$\FP\int_0^1\psi(x)\,\dd x=0$, as follows directly from
$\int_\varepsilon^1\psi(x)\,\dd x=-\log\Gamma(\varepsilon)$.

Here is a short independent check of the lower arithmetic terms.
The Laurent coefficients at degrees two and three give
$\mathsf J_2=0$, $\mathsf J_3=3\zeta(2)$, and
$\delta_2=\delta_3=0$. At degrees four and five,
\[
 \mathsf J_4=12Z_{2,1}=12\zeta(3),\qquad
 \mathsf J_5=60Z_{2,1,1}+10Z_{2,2}
                 +10\zeta(2)^2-\frac{35}{6}\zeta(4).
\]
The classical identities $Z_{2,1,1}=\zeta(4)$,
$Z_{2,2}=3\zeta(4)/4$, and
$\zeta(2)^2=5\zeta(4)/2$ give
$\mathsf J_5=260\zeta(4)/3$.
In the tail notation \eqref{ac:tails},
\[
 \delta_4=-4A_3,\qquad
 \delta_5=5A_2^2-20\zeta(2)A_2-20hA_3+25A_4.
\]
For the consecutive difference use $h=H_{n-1}$,
$A_r=A_r(n)$, and $v=1/n$. These two defects give respectively
\[
 4v^3,\qquad
 -10A_2v^2-20A_3v+20hv^3-30v^4+20\zeta(2)v^2.
\]
Multiplication by $-\log n$ and summation proves
$\mathsf D_4=4\zeta'(3)$ and \eqref{ac:lowerD5} directly from
the convergent strict multiple sums. Equations
\eqref{ac:rawclosed}--\eqref{ac:lowerD5}, together with
\eqref{ac:sixthmain} and \eqref{ac:seventhmain}, therefore give
explicit finite reductions of $Q_6$ and $Q_7$ using only coordinates
defined in this article.

\subsection{Proof ledger for the two higher-degree identities}

For reproducibility, the residue defects that generate every
directional term above are printed in full. In the next two formulas
$h=H_n$ and $A_r=A_r(n)$ as in \eqref{ac:tails}:
\begin{align}\label{ac:delta6}
 \delta_6={}&30A_2^2h-90A_2A_3-120\zeta(2)A_2h
       -60A_3h^2+120\zeta(2)A_3+150A_4h-138A_5,\\
 \delta_7={}&70A_2^3+105A_2^2h^2-315\zeta(2)A_2^2
       -630A_2A_3h+588A_2A_4\notag\\
 &-420\zeta(2)A_2h^2-126\zeta(4)A_2+259A_3^2
       -140A_3h^3+840\zeta(2)A_3h-280\zeta(3)A_3\notag\\
 &+525A_4h^2-651\zeta(2)A_4-966A_5h+833A_6.
       \label{ac:delta7}
\end{align}
They follow by finite coefficient extraction in \eqref{ac:local},
\eqref{ac:appell}, and the elementary harmonic generator. Applying
\eqref{ac:delta-rule} and \eqref{ac:marked-rule} to these two
polynomials gives the unsimplified directional formulas recorded
in the attached exact coefficient tables.

Only the following free-order identities are needed to simplify
their symmetric derivative parts:
\begin{align}
 \zeta(a)\zeta(b)&=Z(a,b)+Z(b,a)+\zeta(a+b),\label{ac:free-stuffle}\\
 Z(s,s,s)&=\frac{\zeta(s)^3-3\zeta(s)\zeta(2s)+2\zeta(3s)}6.
       \label{ac:equal-triple}
\end{align}
Both hold on open absolute-convergence domains. In particular
\begin{align*}
 (\partial_1+\partial_2)Z(2,2)
       &=\zeta(2)\zeta'(2)-\zeta'(4),\\
 (\partial_1+\partial_2)Z(3,3)
       &=\zeta(3)\zeta'(3)-\zeta'(6),\\
 (\partial_1+\partial_2+\partial_3)Z(2,2,2)
       &=\frac34\zeta(4)\zeta'(2)-\zeta(2)\zeta'(4)+\zeta'(6).
\end{align*}
The derivative identities for $(2,3)$ and $(2,4)$ are obtained by
differentiating \eqref{ac:free-stuffle} in either argument. Direct
substitution yields \eqref{ac:D6} and \eqref{ac:D7}. Integer-index
duality is used only for ordinary values in $\mathsf J_m$, never
for these derivatives.

For completeness, the ordinary word calculations before the last
reductions are
\begin{align}\label{ac:J6ledger}
 \mathsf J_6={}&90\zeta(2)\zeta(3)+60\zeta(2)Z_{2,1}
       +360Z_{2,1,1,1}+60Z_{2,1,2}+60Z_{2,2,1}\notag\\
 &-30Z_{2,3}-15\zeta(2)\zeta(3)-15Z_{3,2}
       -35Z_{4,1}-\frac52\zeta(5),\\
 \mathsf J_7={}&588\zeta(4)\zeta(2)+630\zeta(3)Z_{2,1}
       +420\zeta(2)Z_{2,1,1}+2520Z_{2,1,1,1,1}\notag\\
 &+420(Z_{2,1,1,2}+Z_{2,1,2,1}+Z_{2,2,1,1})
       -210(Z_{2,1,3}+Z_{2,3,1})-42Z_{2,4}\notag\\
 &-42\zeta(3)^2-105\zeta(2)Z_{3,1}
       -105(Z_{3,1,2}+Z_{3,2,1})-63Z_{3,3}\notag\\
 &-77\zeta(2)\zeta(4)-245Z_{4,1,1}-77Z_{4,2}
       -\frac{35}{2}Z_{5,1}+\frac{189}{10}\zeta(6).
       \label{ac:J7ledger}
\end{align}
The familiar weight-five Euler and duality identities reduce
\eqref{ac:J6ledger} to $245\zeta(2)\zeta(3)+175\zeta(5)$.
The weight-six entries needed for \eqref{ac:J7ledger} are as
follows. Write $a=\zeta(3)^2$, $b=\zeta(6)$.
\begin{center}
\small
\begin{tabular}{ll}
\hline
Admissible value or equal pair & Reduction\\
\hline
$Z_{2,1,1,1,1}$ & $b$\\
$Z_{2,1,1,2}=Z_{2,4}$ & $25b/12-a$\\
$Z_{2,1,2,1}=Z_{3,3}$ & $(a-b)/2$\\
$Z_{2,2,1,1}=Z_{4,2}$ & $a-4b/3$\\
$Z_{2,1,3}$ & $a-13b/16$\\
$Z_{2,3,1}=Z_{3,1,2}$ & $53b/24-3a/2$\\
$Z_{3,2,1}$ & $3a-203b/48$\\
$Z_{4,1,1}$ & $23b/16-a$\\
$Z_{5,1}$ & $3b/4-a/2$\\
\hline
\end{tabular}
\end{center}
These are ordinary classical multiple-zeta reductions, not new
claims of special-value priority. One finite derivation uses the
strict stuffle product, the binary iterated-integral shuffle product,
and reversal-complement duality \cite{IKZ}. At weight six,
take the products of admissible words of weights $2+4$ and $3+3$,
their reversals, dualities, and the regularized shuffle-minus-stuffle
relation of $(1)$ with each admissible weight-five word. Along with
$Z(6)=\zeta(6)$, this gives 49 explicitly generated equations in
16 admissible weight-six words. The provided rational certificate
contains 16 independent rows and their inverse; multiplication
checks the resulting table exactly. Substitution in
\eqref{ac:J7ledger}, using
$\zeta(2)\zeta(4)=7\zeta(6)/4$, gives the ordinary constant in
\eqref{ac:seventhmain}. This also specifies a short independent
way of auditing its less transparent rational numerator.

\subsection{Validation and the remaining arithmetic question}

The accompanying program constructs the pole coefficients, the
tail polynomials, their finite differences, and the distinguished-index
word maps using exact rational algebra. It checks the weight-six
reduction certificate and the directional simplifications separately.
The analytic statements above justify all infinite operations; a
finite symbolic replay is not a substitute for those statements.

An independent numerical diagnostic compares the locally subtracted
integral of the explicit polynomial \eqref{ac:P6} or
\eqref{ac:P7} with a different representation. In the latter,
$m!(c_m-\lambda_{m-1})$ is evaluated through the ordinary
Gamma-quotient beta kernel established in \cite{Fifth}:
\begin{equation}\label{ac:diagnostic-kernel}
 m!(c_m-\lambda_{m-1})
 =-\int_0^1(-\log y)^m
  \frac{-\log(1-y)-y}{(1-y)^2[-\log(1-y)]^2}\,\dd y.
\end{equation}
The correction \eqref{ac:D} is evaluated from its tail polynomial
at $n\le96$, with the remaining ordinary sum accelerated by
42 inverse-power terms and Hurwitz-zeta order derivatives.
At 80 working decimal digits the comparisons were
\begin{center}
\small
\begin{tabular}{lll}
\hline
Moment & Direct numerical value & Absolute discrepancy\\
\hline
$\mathsf S_6$ & $20.8226131542578224267340551051333\ldots$
 & $8.80\times10^{-56}$\\
$\mathsf S_7$ & $123.619906884273853180167345465232\ldots$
 & $5.73\times10^{-54}$\\
\hline
\end{tabular}
\end{center}
These are floating-point diagnostics, not interval certificates.
They test the finite-part endpoint terms and the ordinary constant
reductions through analytically different formulas.

Theorems~\ref{ac:main}--\ref{ac:seventh} remove an ambiguity from
the source research question. A single singular Laurent coordinate
is sufficient in every degree, while additional \emph{convergent}
derivative coordinates already occur at low degree. The next
arithmetic problem is to determine relations among the explicitly
retained directions in \eqref{ac:D6} and \eqref{ac:D7}, within a
stated relation family or a stated algebra of periods. The depth
bound is a constructive upper bound; it is not a proof that these
directions are necessary or independent.

% END sections/06_appell.tex


% BEGIN sections/07_lerch.tex
\section{Every order derivative on the odd Lerch lattice}
\label{lj:sec:main}

The report \emph{Continuous Depth Lerch} asks in its final question Q8
for the spectral derivatives of its continuous-depth family at odd
positive integers \cite{Continuous}. Its formula \texttt{cdl:eq:oddF} gives the values there,
but does not permit differentiation in the discrete index. We resolve
that question in the precise form requested: an explicit finite formula
for every order derivative and every mixed depth derivative, including
the reflected spectral coordinates. We also give normally convergent
series for all those coordinates. Thus the answer includes a constructive
evaluation procedure at the negative reflected orders, where the original
alternating series need not converge.

The analytic inputs are the source report's proved centered reflection
formula \texttt{cdl:eq:completed}, Euler's classical transformation of
series, and the usual Gamma and Hurwitz-zeta identities \cite{DLMFEuler,DLMFGamma,DLMFZeta}. The new claim is
the following continuation of the inspected source question, not priority
for those classical inputs or for the special cases recovered below.
Euler's transformation is recorded in DLMF \S3.9(ii); the normally
convergent, parameter-differentiated use needed here is proved explicitly.

\subsection{Functions, branches, and normalization}

For $\re\rho>0$, put
\begin{align}
 f_\rho(t)&=(2\cosh(t/2))^{-\rho},\qquad t\geq0,
 \label{lj:eq:f}\\
 Z_\rho(s)&=\frac1{\Gamma(s)}
       \int_0^\infty t^{s-1}f_\rho(t)\,\dd t,
       \qquad \re s>0.
 \label{lj:eq:Z}
\end{align}
The logarithm in $f_\rho$ is real on the positive real axis. The function
$Z_\rho(s)$ is entire in $s$: subtracting finitely many Taylor terms of
$f_\rho$ near zero gives its meromorphic Mellin transform, and the zeros
of $1/\Gamma(s)$ cancel its poles. The same subtractions on compact
parameter sets establish joint holomorphy for $\re\rho>0$. In the
absolutely convergent region $\re s>\re\rho$,
\begin{equation}
 Z_\rho(s)=\sum_{k=0}^\infty
     \frac{(-1)^k(\rho)_k}{k!(k+\rho/2)^s}.
 \label{lj:eq:Zseries}
\end{equation}
Powers of $k+\rho/2$ use the logarithm on the right half-plane.

Define
\begin{align}
 b_\rho(x)&=\frac{(2\pi)^{\rho-2}}{\Gamma(\rho)}
       \Gamma\left(\frac\rho2+\frac{ix}{2\pi}\right)
       \Gamma\left(\frac\rho2-\frac{ix}{2\pi}\right),
 \label{lj:eq:b}\\
 F_\rho(s)&=\frac1{\Gamma(s)}
       \int_0^\infty x^{s-1}b_\rho(x)\,\dd x,
       \qquad \re s>0.
 \label{lj:eq:F}
\end{align}
This is $\mathcal F_{\rho,1/2}(s;0)$ in the continuous-depth report.
Its entire continuation follows from the same Mellin-subtraction argument.
For complex $\rho$, the Gamma product is a product of two functions,
not a squared modulus.

The centered reflection identity is
\begin{equation}
 F_\rho(s)=\frac{\pi(2\pi)^{\rho+s-2}}
                 {\Gamma(s)\sin(\pi s/2)}Z_\rho(1-s).
 \label{lj:eq:reflection}
\end{equation}
Apparent singularities are removed in the combined expression. For
completeness, the beta integral gives the classical Fourier formula \cite{RamanujanFourier}
$\widehat f_\rho(v)=
\Gamma(\rho/2+iv)\Gamma(\rho/2-iv)/\Gamma(\rho)$, where the Fourier
transform is $\int_{\mathbb R}e^{ivt}f_\rho(t)\,\dd t$. On
$0<\re s<1$, insert a decaying exponential in the cosine Mellin
integral, use Fubini, and then dominated convergence to obtain
\[
 \int_0^\infty v^{s-1}\widehat f_\rho(v)\,\dd v
 =2\Gamma(s)\cos(\pi s/2)
       \int_0^\infty t^{-s}f_\rho(t)\,\dd t.
\]
The regulator can be removed on the left because the Gamma product
decays exponentially. Rescaling gives
$F_\rho(s)=2(2\pi)^{\rho+s-2}\Gamma(1-s)
\cos(\pi s/2)Z_\rho(1-s)$; Euler reflection for Gamma gives
\eqref{lj:eq:reflection}. Entire continuation gives the stated global
interpretation. This reproduces the source input rather than replacing
it with a formal exchange at a pole.

We use $\partial_s^q$ for order derivatives and $\partial_\rho^r$ for
depth derivatives. The complete exponential Bell polynomials are fixed by
\begin{equation}
 \exp\left(\sum_{j\geq1}d_j\frac{v^j}{j!}\right)
 =\sum_{q\geq0}\mathcal B_q(d_1,\ldots,d_q)\frac{v^q}{q!},
 \qquad \mathcal B_0=1.
 \label{lj:eq:Bell}
\end{equation}
Write $H_N^{(j)}=\sum_{k=1}^N k^{-j}$ and $H_0^{(j)}=0$.

\subsection{The finite formula for all mixed jets}

\begin{theorem}[Odd-lattice order and depth derivatives]
\label{lj:thm:oddjets}
Let $m,r,q\geq0$ be integers, put $N=2m$, $n=N+1$, and set
\begin{align}
 L&=\log(2\pi),\qquad
 C_{m,\rho}=\frac{(-1)^m\pi(2\pi)^{\rho+N-1}}{N!},
 \label{lj:eq:C}\\
 \kappa_1&=L+\gamma-H_N,\nonumber\\
 \kappa_j&=(j-1)!\left[
       \left(1-\frac{1+(-1)^j}{2^j}\right)\zeta(j)
       +(-1)^jH_N^{(j)}\right],\qquad j\geq2.
 \label{lj:eq:kappa}
\end{align}
For every $\re\rho>0$,
\begin{equation}
 \boxed{\displaystyle
 \partial_\rho^r\partial_s^qF_\rho(n)
 =C_{m,\rho}\sum_{j=0}^q(-1)^j\binom qj
       \mathcal B_{q-j}(\kappa_1,\ldots,\kappa_{q-j})
       (\partial_\rho+L)^r\partial_s^jZ_\rho(-N).}
 \label{lj:eq:alljets}
\end{equation}
The operator on the right is explicitly
$\sum_{h=0}^r\binom rh L^{r-h}\partial_\rho^h$.
No differentiation with respect to $m$ occurs.
\end{theorem}

\begin{proof}
The pole-free local identity from which every derivative follows is
\begin{equation}
 F_{\rho+u}(n+v)=C_{m,\rho}e^{Lu}
       R_m(v)Z_{\rho+u}(-N-v),
 \qquad
 R_m(v)=e^{Lv}\frac{\Gamma(n)}{\Gamma(n+v)}
                  \sec(\pi v/2).
 \label{lj:eq:local}
\end{equation}
It holds for $u$ sufficiently small and $|v|<1$, by substituting in
\eqref{lj:eq:reflection} and using
$\sin(\pi(n+v)/2)=(-1)^m\cos(\pi v/2)$.
All factors in \eqref{lj:eq:local} are holomorphic there.

The Gamma contribution to the $j$th logarithmic derivative of $R_m$
is $-\psi^{(j-1)}(n)$. For $j\geq2$,
\[
 \psi^{(j-1)}(n)=(-1)^j(j-1)!
                         \bigl(\zeta(j)-H_N^{(j)}\bigr).
\]
The product for cosine gives the convergent expansion
\[
 \log\sec(\pi v/2)
   =\sum_{k\geq1}\frac{(1-2^{-2k})\zeta(2k)}{k}v^{2k},
   \qquad |v|<1.
\]
Combining these expansions gives exactly \eqref{lj:eq:kappa} and
$R_m^{(a)}(0)=\mathcal B_a(\kappa_1,\ldots,\kappa_a)$.
Now apply the ordinary multivariate Leibniz rule to the holomorphic
identity \eqref{lj:eq:local}. The argument $-N-v$ supplies $(-1)^j$,
while $e^{Lu}$ supplies $(\partial_\rho+L)^r$. This proves
\eqref{lj:eq:alljets} with every combinatorial coefficient fixed.
\end{proof}

For example, putting $Z^{[j]}=\partial_s^jZ_\rho(-N)$, the first two
order derivatives are
\begin{align}
 F_\rho'(n)&=C_{m,\rho}(\kappa_1Z^{[0]}-Z^{[1]}),
 \label{lj:eq:first}\\
 F_\rho''(n)&=C_{m,\rho}
       \bigl((\kappa_1^2+\kappa_2)Z^{[0]}
                      -2\kappa_1Z^{[1]}+Z^{[2]}\bigr),
 \label{lj:eq:second}\\
 \kappa_2&=\tfrac12\zeta(2)+H_N^{(2)},\qquad
 \kappa_3=2\bigl(\zeta(3)-H_N^{(3)}\bigr).\nonumber
\end{align}

To recover the already known value, define
$U_m(\rho)=[z^{2m}]\sec(z)^\rho$. The even Taylor expansion of
$f_\rho$ gives
\begin{equation}
 Z_\rho(-2m)=(-1)^m(2m)!2^{-\rho-2m}U_m(\rho),
 \qquad
 F_\rho(2m+1)=\frac{\pi^{\rho+2m}}2U_m(\rho).
 \label{lj:eq:valuecheck}
\end{equation}
These are the $q=0$ inputs. The terms $Z^{[j]}$ with $j>0$ in
\eqref{lj:eq:alljets} are additional data; discarding them would give
an incorrect spectral derivative.

\subsection{A convergent elementary series for every reflected jet}

The next theorem makes the reflected coordinates in
\eqref{lj:eq:alljets} independently computable. In particular, it avoids
using the reflection equation itself as a numerical definition of the
coordinates that are supposed to verify that equation.

\begin{theorem}[Normally convergent Euler series at every order]
\label{lj:thm:Euler}
For $\re\rho>0$ and $s\in\mathbb C$, define
\begin{equation}
 A_k(\rho,s)=\frac{(\rho)_k}{k!(k+\rho/2)^s},\qquad
 D_\ell(\rho,s)=\sum_{k=0}^\ell(-1)^k\binom\ell k A_k(\rho,s).
 \label{lj:eq:D}
\end{equation}
Then
\begin{equation}
 \boxed{\displaystyle
 Z_\rho(s)=\sum_{\ell=0}^\infty\frac{D_\ell(\rho,s)}{2^{\ell+1}}.}
 \label{lj:eq:Euler}
\end{equation}
This series is absolutely and locally normally convergent throughout
$\{\re\rho>0\}\times\mathbb C$. Every fixed mixed derivative in
$\rho,s$ is obtained by differentiating its finite inner sums termwise.
\end{theorem}

\begin{proof}
For $z\in\Omega=\mathbb C\setminus[1,\infty)$, initially $\re s>0$,
let
\begin{equation}
 A(\rho,s;z)=\frac1{\Gamma(s)}\int_0^\infty
        t^{s-1}e^{-\rho t/2}(1-ze^{-t})^{-\rho}\,\dd t.
 \label{lj:eq:Aint}
\end{equation}
Use the principal logarithm of $1-ze^{-t}$. It is available on
$\Omega$, since $1-ze^{-t}$ cannot lie on the nonpositive real axis
there. On compact subsets, the integrand is bounded at zero after the
factor $t^{s-1}$ is removed and decays exponentially at infinity.
Thus \eqref{lj:eq:Aint} is jointly holomorphic in this initial domain.
Its power series at $z=0$ is $\sum_{k\geq0}A_k(\rho,s)z^k$.

The coefficient identity
\begin{equation}
 \left(z\frac\partial{\partial z}+\frac\rho2\right)
 A(\rho,s+1;z)=A(\rho,s;z)
 \label{lj:eq:Aladder}
\end{equation}
holds first in $|z|<1$ and then throughout $\Omega$. Iterating it from
$\re(s+M)>0$ continues $A$ to every $s\in\mathbb C$. The definitions
agree on overlaps by the same identity, so this is a joint holomorphic
continuation. At $z=-1$, it equals $Z_\rho(s)$.

Set
\[
 B(\rho,s;v)=\frac1{1-v}
              A\left(\rho,s;\frac{-v}{1-v}\right).
\]
If $|v|<1$, then $-v/(1-v)$ is outside $[1,\infty)$: a point of that
cut would require $v\in(1,\infty)$. Hence $B$ is holomorphic for
$|v|<1$. Its Taylor coefficient is $D_\ell(\rho,s)$, by expanding
$(1-v)^{-k-1}$ in the defining power series near zero.
Taking $v=1/2$ yields $B(\rho,s;1/2)=2Z_\rho(s)$ and proves
\eqref{lj:eq:Euler}.

For a compact parameter set, choose any fixed radius $r$ with
$1/2<r<1$. Cauchy's coefficient estimate on $|v|=r$ gives
$|D_\ell|\leq C r^{-\ell}$ uniformly on that set. It proves absolute
normal convergence of \eqref{lj:eq:Euler}. Apply the same argument on
a slightly enlarged compact parameter neighborhood to justify every
fixed mixed derivative. This argument proves convergence even when
\eqref{lj:eq:Zseries} is a divergent ordinary series.
\end{proof}

All finite derivative coefficients admit a harmonic expression. Put
$a_k=k+\rho/2$ and
\[
 h_k^{(j)}(\rho)=\sum_{\ell=0}^{k-1}(\rho+\ell)^{-j},\qquad
 d_j=(-1)^{j-1}(j-1)!
       \left(h_k^{(j)}(\rho)-\frac{s}{2^ja_k^j}\right).
\]
Then $\partial_\rho^r A_k=A_k\mathcal B_r(d_1,\ldots,d_r)$.
For mixed derivatives a particularly safe finite generating formula is
\begin{equation}
 \frac{A_k(\rho+u,s+v)}{A_k(\rho,s)}
 =e^{-v\log a_k}
       \prod_{\ell=0}^{k-1}\left(1+\frac{u}{\rho+\ell}\right)
       \left(1+\frac{u}{2a_k}\right)^{-s-v}.
 \label{lj:eq:mixedcoeff}
\end{equation}
Only the requested bidegree is needed. This also distinguishes a mixed
derivative from the product of two independent logarithmic derivatives.

At the reflected integer $s=-N$, $N\geq0$, the formula is especially
explicit:
\begin{equation}
\begin{split}
 \partial_\rho^r\partial_s^q A_k(\rho,-N)
 &=(-1)^q r!\frac{(\rho)_k}{k!}a_k^N[u^r]
       \prod_{\ell=0}^{k-1}\left(1+\frac{u}{\rho+\ell}\right)
       \left(1+\frac{u}{2a_k}\right)^N\\
 &\hspace{13mm}\cdot
       \left(\log a_k+\log\left(1+\frac{u}{2a_k}\right)\right)^q.
\end{split}
 \label{lj:eq:negativecoeff}
\end{equation}
For positive rational $\rho$, every inner coefficient here is a rational
linear combination of powers of logarithms of positive rational numbers.
Equations \eqref{lj:eq:negativecoeff}, \eqref{lj:eq:Euler}, and
\eqref{lj:eq:alljets} therefore constitute a complete ordinary convergent
series evaluation of all the reflected mixed jets in the theorem.
The outer Euler series need not terminate. At a positive integer $\rho$
and $q=r=0$, its inner sequence is polynomial at negative integer $s$,
so finite differences do terminate; differentiation can remove that
termination.

\subsection{Independent Gamma-moment realization}

The positive-half-line Gamma integral gives a different exact
representation for checking the lattice formulas. Put
\begin{equation}
 a_j(\rho,x)=\delta_{j1}\log(2\pi)-\psi^{(j-1)}(\rho)
  +2^{-j}\left\{
   \psi^{(j-1)}\left(\frac\rho2+\frac{ix}{2\pi}\right)
  +\psi^{(j-1)}\left(\frac\rho2-\frac{ix}{2\pi}\right)\right\}.
 \label{lj:eq:Gammajet}
\end{equation}
For the order derivatives define $c_1(x;s)=\log x-\psi(s)$ and
$c_j(x;s)=-\psi^{(j-1)}(s)$ for $j\geq2$. Direct differentiation of
\eqref{lj:eq:F} gives, for $\re s>0$,
\begin{equation}
\begin{split}
 \partial_\rho^r\partial_s^qF_\rho(s)
  =\frac1{\Gamma(s)}\int_0^\infty
      x^{s-1}b_\rho(x)
      \mathcal B_r(a_1(\rho,x),\ldots,a_r(\rho,x))\qquad\qquad\\[-2pt]
      {}\cdot\mathcal B_q(c_1(x;s),\ldots,c_q(x;s))\,\dd x.
\end{split}
 \label{lj:eq:momentjets}
\end{equation}
This is an ordinary absolutely convergent integral. Near zero, only
fixed powers of $\log x$ occur. At infinity, the Gamma product decays
as $e^{-x/2}$ times a power of $x$; its parameter derivatives introduce
only logarithmic or polynomial factors on compact parameter sets.
These observations justify every differentiation. The left side of
\eqref{lj:eq:alljets} can therefore be evaluated through
\eqref{lj:eq:momentjets} while its right side is evaluated through
the independent Euler series.

\subsection{Stieltjes coordinates at depth two}

Let the generalized Stieltjes constants have the convention
\[
 \zeta(1+v,a)=\frac1v+
       \sum_{j\geq0}\frac{(-1)^j\gamma_j(a)}{j!}v^j,
       \qquad \re a>0.
\]
An upper parenthesized index on $\gamma_j$ below is an argument
derivative. Set $E_a(s)=s\zeta(s+1,a)$, with the removable value
$E_a(0)=1$.

\begin{proposition}[Stieltjes transport to every positive order]
\label{lj:prop:Stieltjes}
For integers $n\geq1,q\geq0$, put $N=n-1$ and
$\lambda_j=(-1)^j(j-1)!H_N^{(j)}$. Then
\begin{equation}
 \boxed{\displaystyle
 E_a^{[q]}(n)=\frac{(-1)^n}{N!}
   \sum_{j=0}^q(-1)^j\binom qj
       \mathcal B_{q-j}(\lambda_1,\ldots,\lambda_{q-j})
       \gamma_j^{(n)}(a).}
 \label{lj:eq:Stieltjes}
\end{equation}
In particular,
\begin{equation}
 E_a^{[q]}(1)=(-1)^{q+1}\gamma_q'(a).
 \label{lj:eq:Stieltjesone}
\end{equation}
At $a=1/2$, this gives $\partial_s^qF_2(n)$.
\end{proposition}

\begin{proof}
The shift derivative of Hurwitz zeta gives
$\partial_a^n\zeta(1+v,a)=(-1)^n(1+v)_n\zeta(n+1+v,a)$.
The pole $1/v$ is independent of $a$, so
\begin{equation}
 E_a(n+v)=\frac{(-1)^n}{(1+v)_N}
       \sum_{j\geq0}\frac{(-1)^j\gamma_j^{(n)}(a)}{j!}v^j.
 \label{lj:eq:Stieltjesgerm}
\end{equation}
The logarithmic Taylor coefficients of $N!/(1+v)_N$ are $\lambda_j$.
Coefficient extraction proves \eqref{lj:eq:Stieltjes}.
Finally, $b_2(x)=x/(2\sinh(x/2))$, so its Mellin transform gives
$F_2(s)=s\zeta(s+1,1/2)$ directly. This proves the last assertion.
\end{proof}

This proposition is an explicit form of the familiar parameter-derivative
tower. Its role here is to transport the new lattice-jet formula into
the repository's Stieltjes convention; it is not presented as a new
Hurwitz-zeta identity. For instance, with $N=2m$,
\[
 F_2'(2m+1)=\frac{
       \gamma_1^{(2m+1)}(1/2)
        +H_{2m}\gamma_0^{(2m+1)}(1/2)}{(2m)!}.
\]
Together with \eqref{lj:eq:first}, this is an exact relation between
negative reflected $Z_2$ order derivatives and positive-order Stieltjes
argument derivatives.

\subsection{Gamma-weighted polylogarithm order derivatives}

Put
\begin{align}
 G_\rho(x)&=
   \Gamma\left(\frac\rho2+\frac{ix}{2\pi}\right)
   \Gamma\left(\frac\rho2-\frac{ix}{2\pi}\right),\nonumber\\
 P_s(x)&=\frac{\operatorname{Li}_s(ie^{-x/2})
                    -\operatorname{Li}_s(-ie^{-x/2})}{i},\nonumber\\
 I_\rho(s)&=\int_{\mathbb R}G_\rho(x)P_s(x)\,\dd x.
 \label{lj:eq:I}
\end{align}
The polylogarithms use their principal branches. Their arguments stay
off the positive real cut. The following formula covers all independent
order derivatives, including orders for which the integrand is no longer
an elementary logarithm or a classical fixed-order polylogarithm.

\begin{theorem}[All polylogarithmic order jets under a Gamma weight]
\label{lj:thm:polylog}
For $\re\rho>0$, the integral $I_\rho(s)$ is entire in $s$, jointly
holomorphic in its parameters, and
\begin{equation}
 I_\rho(s)=\Gamma(\rho)(2\pi)^{2-\rho}2^{1-s}F_{\rho+1}(s)
 =\frac{4\Gamma(\rho)\pi^{s+2}}
        {\Gamma(s)\sin(\pi s/2)}Z_{\rho+1}(1-s),
 \label{lj:eq:Imaster}
\end{equation}
with all apparent poles removed in the combined expression.
For $N=2m$, $n=N+1$, replace only the first cumulant in
\eqref{lj:eq:kappa} by
$\widetilde\kappa_1=\log\pi+\gamma-H_N$, and leave
$\widetilde\kappa_j=\kappa_j$ for $j\geq2$. Then
\begin{equation}
 \partial_s^qI_\rho(n)=
 \frac{4(-1)^m\Gamma(\rho)\pi^{n+2}}{N!}
   \sum_{j=0}^q(-1)^j\binom qj
    \mathcal B_{q-j}(\widetilde\kappa_1,\ldots,
                    \widetilde\kappa_{q-j})
    \partial_s^jZ_{\rho+1}(-N).
 \label{lj:eq:Ijets}
\end{equation}
Every mixed depth derivative follows by the finite rule
\begin{equation}
 \partial_\rho^r\{\Gamma(\rho)Y(\rho+1)\}
 =\Gamma(\rho)\sum_{h=0}^r\binom rh
     \mathcal B_{r-h}(\psi(\rho),\psi'(\rho),\ldots)
     Y^{(h)}(\rho+1).
 \label{lj:eq:Igammajets}
\end{equation}
All these derivatives pass through the ordinary integral in
\eqref{lj:eq:I}.
\end{theorem}

\begin{proof}
The root filter in the continuous-depth report, valid first by a
convergent series and then by analytic continuation, is
\[
 \mathcal F_{1,1/2}(s;-x)=2^{s-1}e^{-x/2}P_s(x).
\]
Convolve this with
$K_\rho(x)=(2\pi)^{\rho-2}e^{x/2}G_\rho(x)/\Gamma(\rho)$.
The source's entire-order addition law gives $F_{\rho+1}(s)$ and
therefore the first equality in \eqref{lj:eq:Imaster}. The second is
\eqref{lj:eq:reflection}, with the powers of $2$ and $\pi$ collected.
The ordinary integral is absolutely convergent locally uniformly in
all the stated parameters. One explicit justification uses the vertical
contour for $\mathcal F_{1,1/2}$ on any line $1/2<c<1$, giving
$|P_s(x)|\leq C e^{(1/2-c)x}$ on compact order sets. The Gamma weight
then decays integrably at both ends. Fixed order derivatives add powers
of the logarithm of the contour variable, and depth derivatives of the
Gamma weight add only polynomial or logarithmic factors. The same
integrable majorants persist on slightly larger compact sets, proving
joint holomorphy and passage of all derivatives.

The local germ at $s=n+v$ now contains $e^{v\log\pi}$ in place of
$e^{v\log(2\pi)}$. The proof of Theorem~\ref{lj:thm:oddjets} gives
\eqref{lj:eq:Ijets}. Finally, the logarithmic derivatives of
$\Gamma(\rho)$ are $\psi^{(j-1)}(\rho)$, so Bell's formula and Leibniz
give \eqref{lj:eq:Igammajets}.
\end{proof}

On the integral side of \eqref{lj:eq:Igammajets}, the weight replacing
$G_\rho$ is
\[
 G_\rho(x)\mathcal B_r(g_1(\rho,x),\ldots,g_r(\rho,x)),\qquad
 g_j(\rho,x)=2^{-j}\left\{
 \psi^{(j-1)}\left(\frac\rho2+\frac{ix}{2\pi}\right)
 +\psi^{(j-1)}\left(\frac\rho2-\frac{ix}{2\pi}\right)\right\}.
\]
Thus the theorem is a family of explicit ordinary integrals containing
polygamma factors and spectral derivatives of polylogarithms. Theorem
\ref{lj:thm:Euler} evaluates their reflected coordinates by convergent
series, including the mixed derivatives appearing in
\eqref{lj:eq:Igammajets}.

At $\rho=1$ and $s=1$, $G_1(x)=\pi\operatorname{sech}(x/2)$.
The first derivative has the concrete forms
\begin{align}
 \int_{\mathbb R}\frac{\pi}{\cosh(x/2)}
       \left.\partial_sP_s(x)\right|_{s=1}\,\dd x
 &=2\pi\gamma_1'(1/2)-\pi^3\log2
 \label{lj:eq:exampleStieltjes}\\
 &=\pi^3\left(\gamma+\log\pi+
                  \frac43\log2+12\zeta'(-1)\right).
 \label{lj:eq:exampleZeta}
\end{align}
Indeed, the first equality follows from
$I_1(s)=2\pi2^{1-s}E_{1/2}(s)$ and
\eqref{lj:eq:Stieltjesone}. For the second use
$Z_2(s)=\eta(s-1)$,
$\eta(s)=(1-2^{1-s})\zeta(s)$,
$Z_2(0)=1/4$, and
$Z_2'(0)=-\log2/3-3\zeta'(-1)$ in
\eqref{lj:eq:Ijets}. The second order derivative is
\begin{equation}
 \int_{\mathbb R}\frac{\pi}{\cosh(x/2)}
       \left.\partial_s^2P_s(x)\right|_{s=1}\,\dd x
 =-2\pi\gamma_2'(1/2)-4\pi\log2\,\gamma_1'(1/2)
       +\pi^3\log^22.
 \label{lj:eq:exampleSecond}
\end{equation}
These formulas state the normalization explicitly; in particular, the
$2^{1-s}$ factor cannot be held constant during order differentiation.

\subsection{What this resolves and what remains}

Theorem~\ref{lj:thm:oddjets} resolves Q8 in the source's requested
Gamma--Stieltjes--$Z_\rho$-jet form. Theorem~\ref{lj:thm:Euler} supplies
a convergent evaluator and explicit harmonic coefficient calculus for
every reflected coordinate, while Theorem~\ref{lj:thm:polylog} transports
the result to ordinary Gamma-weighted polylogarithm integrals. This does
not assert that arbitrary noninteger-depth spectral jets reduce to a
finite algebra of classical zeta and Gamma values. Such a further
reduction would be a different arithmetic statement.

No false theorem was found in the inspected continuous-depth report.
Its caution that one cannot differentiate a formula known only at
integer $m$ is correct. The necessary addition is a holomorphic local
identity, \eqref{lj:eq:local}, together with the extra reflected jets.
The tilted report already treats independent shift and order derivatives
of its matched convolution family; those results are retained as prior
work and are not recounted as discoveries here.

Two concrete questions remain. First, determine whether selected rational
noninteger depths admit additional finite reductions of
$\partial_s^qZ_\rho(-2m)$, specifying the comparison algebra in advance.
Second, develop a cancellation-aware implementation of the Euler sums
with proved arithmetic error bounds. The analytic theorem supplies
normal convergence, while finite floating-point agreement by itself
does not provide a certificate for a transcendental identity.

% END sections/07_lerch.tex


% BEGIN sections/08_adjacent_tails.tex
% Self-contained research section for integration into the parent article.
% Requires amsmath, amssymb, amsthm; standard theorem environments.
% Source target: ProveIt_Canonical_Digamma_Fifth_Moment_2026-10-11,
% sections/09_questions.tex, Conjecture conj:firstq; its inherited matrix
% normalization is sections/07_tail_rank.tex, equations tr:Aentries,
% tr:column and tr:operator. No source files have been modified.

\section{Optimal initial moments for adjacent even tail orders}
\label{ar:section}

The incoming canonical-digamma report \cite{Fifth} conjectures that the first $q$
centered moments of two unequal even harmonic tails have full formal
rank $q$ (its Conjecture~\texttt{conj:firstq}). That report proves the
claim when the two orders coincide and gives a larger sufficient prefix
for every unequal pair. We prove the conjectured optimal prefix for the
entire next family
\[
                 (p,q)=(2b-2,2b),\qquad b\ge2.
\]
The proof also gives an exact determinant with a strictly positive
rational recurrence. Thus the nonvanishing is an all-order theorem,
independent of numerical experiments. The conjecture for general larger
gaps remains open here.

\subsection{The inherited coordinates and the exact target}

For even integers $2\le p\le q=2b$, put
\[
 A_p(n)=\zeta(p)-H_n^{(p)},\qquad
 D_{p,q}(s)=\sum_{n\ge0}
     A_p(n)A_q(n)(n+\tfrac12)^{-s},\qquad
 D_r=D_{p,q}(-2r).
\]
The series initially defines $D_{p,q}$ for
$\Re s>3-p-q$; the notation $D_r$ uses its meromorphic continuation.
The centered asymptotic series for the product of the two tails contains
only even powers, so every point $s=-2r$ is regular. This regularity and
the following finite coordinate formulas are inherited from the two-tail
calculus, not additional conclusions of the rank argument.

The source supplies blocks $C_j$ with mutually inverse transforms
\begin{align}
 D_r&=\sum_{j=0}^r
    \frac{\binom{2r}{2j}}{2j+1}B_{2r-2j}(\tfrac12)C_j,
       \label{ar:forward}\\
 C_j&=\sum_{r=0}^j
       \binom{2j+1}{2r}2^{-2(j-r)}D_r.
       \label{ar:backward}
\end{align}
Here $B_n(x)$ is the Bernoulli polynomial and $B_n=B_n(0)$, with
$B_1=-1/2$. The first $b$ blocks $C_0,\ldots,C_{b-1}$ are kept as
specified coordinate labels. All subsequent blocks have the form
\begin{equation}\label{ar:block}
 C_j=\sum_{k=0}^{b-1}a_{jk}Z_k,\qquad
 Z_k=\zeta(2k),\quad Z_0=-\tfrac12,\qquad j\ge b.
\end{equation}
For $d=2j+1$ and $e=d-p-q+2k+1$, the exact rational coefficients are
\begin{equation}\label{ar:entries}
 a_{jk}=\begin{cases}
 B_e\left\{
 \displaystyle\frac{\binom{d-q+1}{p-2k}}{d-q+1}
 +\displaystyle\frac{\binom{d-p+1}{q-2k}}{d-p+1}
 \right\},&e\ge0,\\[6pt]
 0,&e<0.
 \end{cases}
\end{equation}
A binomial summand is omitted unless its lower index is at least two
and no larger than its upper index. In particular a putative lower
index zero is omitted. This endpoint convention is essential to the
proof and to the determinant.

We consider the $b\times b$ late-block matrix
\begin{equation}\label{ar:M}
 M_{p,q}=(a_{b+r-1,k})_{1\le r\le b,\ 0\le k<b}.
\end{equation}
The $2b\times2b$ coefficient matrix of
$D_0,\ldots,D_{2b-1}$ on the formal labels
\[
           C_0,\ldots,C_{b-1},Z_0,\ldots,Z_{b-1}
\]
is nonsingular exactly when $M_{p,q}$ is nonsingular. Indeed the
corresponding block matrix for $C_0,\ldots,C_{2b-1}$ is
$\operatorname{diag}(I_b,M_{p,q})$, and the transform
\eqref{ar:forward} is triangular with nonzero diagonal
$1/(2r+1)$. These are formal ranks over $\mathbb Q$; they imply
generation statements for the actual values without asserting their
arithmetic independence.

\subsection{A positive recurrence and the determinant}

Define positive rational numbers
\begin{equation}\label{ar:alpha}
 \alpha_n=\frac{2|B_{2n}|}{(2n)!}
         =\frac{4\zeta(2n)}{(2\pi)^{2n}},\qquad n\ge1.
\end{equation}
The equality follows from Euler's evaluation of even zeta values and
establishes positivity at every index. Define $\rho_1=1$ and, for $r\ge1$,
\begin{equation}\label{ar:rho}
 \rho_{r+1}=
 \frac{1}{4(2r+1)^2}
 \sum_{s=1}^r\alpha_{r+1-s}
  \left\{2r(2r+1)
       +\mathbf1_{s\ge2}\,2s(2s-1)\right\}\rho_s.
\end{equation}
Every coefficient in this recurrence is strictly positive. In particular
$\rho_s>0$ for every $s$. The beginning is
\[
 \rho_1=1,\qquad \rho_2=\frac1{36},\qquad
 \rho_3=\frac{11}{5400},\qquad
 \rho_4=\frac{113}{705600},\qquad
 \rho_5=\frac{2509}{190512000}.
\]

\begin{theorem}[The optimal prefix for adjacent even orders]
\label{ar:main}
For $q=2b\ge4$ and $p=q-2$, the matrix in \eqref{ar:M} has determinant
\begin{equation}\label{ar:det}
 \boxed{\displaystyle
 \det M_{q-2,q}
 =\frac{(-1)^{b(b-1)/2}}{b!}
       \bigl((2b+1)!!\bigr)^2\rho_{b+1}\ne0.}
\end{equation}
Consequently the first $q$ moments $D_0,\ldots,D_{q-1}$ have full
formal rank $q$, and
\begin{equation}\label{ar:spans}
 \operatorname{span}_{\mathbb Q}\{D_0,D_1,\ldots\}
 =\operatorname{span}_{\mathbb Q}\{D_0,\ldots,D_{q-1}\}
 =\operatorname{span}_{\mathbb Q}
     \{C_0,\ldots,C_{b-1},Z_0,\ldots,Z_{b-1}\}.
\end{equation}
The complete first-$q$ moment matrix, in the displayed coordinate order,
has determinant
\begin{equation}\label{ar:complete}
 \frac{(-1)^{b(b-1)/2}\bigl((2b+1)!!\bigr)^2\rho_{b+1}}
      {b!\,(4b-1)!!}.
\end{equation}
No prefix of fewer than $q$ rows can have rank $q$ on these $q$ formal
labels, so the prefix length is optimal in this specified formal system.
\end{theorem}

\subsection{Proof by rotation of the local differential equation}

We reproduce the source's differential description so that the argument
can be checked directly against \eqref{ar:entries}. Put $h(t)=\coth(t/2)$.
For $0\le k<b$, define the analytic germs
\begin{align}
 f_k(t)={}&\frac{1}{2(p-2k)!(q-1)!}
      \int_0^t u^{p-2k}(t-u)^{q-1}h(u)\,du\notag\\
 &+\frac{1}{2(q-2k)!(p-1)!}
      \int_0^t u^{q-2k}(t-u)^{p-1}h(u)\,du.
 \label{ar:column}
\end{align}
The first summand is omitted when $p-2k<2$. All included integrands
are regular at $u=0$. Their power series converge in $|t|<2\pi$.
Expanding
\[
        h(u)=2\sum_{n\ge0}\frac{B_{2n}}{(2n)!}u^{2n-1}
\]
and integrating monomials shows that the coefficient of
$t^{2j+1}/(2j+1)!$ in $f_k$ is precisely $a_{jk}$ for $j\ge b$.

For $v=(v_0,\ldots,v_{b-1})$ set
\[
 f_v=\sum_{k=0}^{b-1}v_kf_k,
 \qquad Q(t)=\sum_{k=0}^{b-1}v_k\frac{t^{q-2k}}{(q-2k)!}.
\]
When $p=q-2$, differentiating \eqref{ar:column} $q$ times gives
\begin{equation}\label{ar:operator}
 2f_v^{(q)}(t)
    =\bigl(Q''(t)-Q''(0)\bigr)h(t)+(Q(t)h(t))''
    =:\mathcal L_Q(t).
\end{equation}
The subtracted constant removes exactly the omitted lower-index-zero
term. Each side is an odd analytic germ at zero.

Rotate $t=ix$, put
\[
 R(x)=Q(ix)=\sum_{s=1}^b r_sx^{2s},\qquad
 r_s=\frac{(-1)^s}{(2s)!}v_{b-s},
 \qquad c(x)=\cot(x/2).
\]
Because $h(ix)=-ic(x)$ and two differentiations contribute a factor
$-1$, \eqref{ar:operator} gives
\begin{equation}\label{ar:rotated}
 \mathcal L_Q(ix)=i\bigl((R''-R''(0))c+(Rc)''\bigr).
\end{equation}
The benefit of this rotation is the one-sided coefficient expansion
\begin{equation}\label{ar:cot}
 c(x)=\frac2x-\sum_{n\ge1}\alpha_nx^{2n-1},\qquad |x|<2\pi,
\end{equation}
where every $\alpha_n$ is positive.

Let $Y_r$ be the ordinary coefficient of $x^{2r-1}$ in the bracket
on the right of \eqref{ar:rotated}. Write $r_{b+1}=0$ because $R$
has degree at most $2b$. Direct differentiation of the monomials in
\eqref{ar:cot} gives, for $1\le r\le b$,
\begin{equation}\label{ar:coefficient}
 Y_r=d_{r+1}r_{r+1}-\sum_{s=1}^r w_{rs}r_s,
\end{equation}
where
\begin{equation}\label{ar:weights}
 d_{r+1}=4(2r+1)^2,\qquad
 w_{rs}=\alpha_{r+1-s}
  \left\{2r(2r+1)+\mathbf1_{s\ge2}\,2s(2s-1)\right\}>0.
\end{equation}
For clarity, the coefficient supplied by $2/x$ and $r_sx^{2s}$,
$s\ge2$, is
\[
 2\{(2s)(2s-1)+(2s-1)(2s-2)\}r_s
       =4(2s-1)^2r_s.
\]
For $s=1$, the first differentiation term is removed by $R''(0)$,
while the second is the second derivative of a linear monomial and
vanishes. This explains both the superdiagonal $d_{r+1}$ and the
exceptional indicator in \eqref{ar:weights}.

Suppose first that $M_{q-2,q}v=0$. The first $b$ odd coefficients
of $\mathcal L_Q$ vanish, hence $Y_1=\cdots=Y_b=0$.
The first $b-1$ equations in \eqref{ar:coefficient} force
\[
                   r_s=\rho_sr_1,\qquad 1\le s\le b.
\]
The last equation, using $r_{b+1}=0$, then gives
\[
           0=-\sum_{s=1}^bw_{bs}\rho_sr_1
             =-d_{b+1}\rho_{b+1}r_1.
\]
All factors preceding $r_1$ are nonzero. Therefore $r_1=0$, then
$R=Q=0$ and $v=0$. This already proves nonsingularity over
$\mathbb C$, and therefore over $\mathbb Q$.

To retain the exact determinant, let $H_b$ be the matrix of
\eqref{ar:coefficient} from $(r_1,\ldots,r_b)$ to
$(Y_1,\ldots,Y_b)$. Replace its first column basis vector by
$(\rho_1,\ldots,\rho_b)^{\mathsf T}$, leaving the others fixed.
The basis change has determinant one. Its first column image has
only the final entry $-d_{b+1}\rho_{b+1}$. The first $b-1$ rows
on columns $2,\ldots,b$ form a lower triangular matrix with diagonal
$d_2,\ldots,d_b$. Expansion along the first column yields
\begin{equation}\label{ar:Hdet}
 \det H_b=(-1)^b\rho_{b+1}\prod_{s=2}^{b+1}d_s
    =(-1)^b4^b\bigl((2b+1)!!\bigr)^2\rho_{b+1}.
\end{equation}
On the other hand, coefficient comparison in \eqref{ar:rotated}
and \eqref{ar:operator} gives
\begin{equation}\label{ar:scalings}
 Y_r=\frac{2(-1)^{r-1}}{(2r-1)!}(M_{q-2,q}v)_r,
 \qquad r_s=\frac{(-1)^s}{(2s)!}v_{b-s}.
\end{equation}
The reversal of the $b$ columns has determinant
$(-1)^{b(b-1)/2}$. Combining the diagonal factors in
\eqref{ar:scalings} gives
\[
 \det M_{q-2,q}
     =\frac{(-1)^{b(b+1)/2}}{4^bb!}\det H_b.
\]
Substitute \eqref{ar:Hdet}; the resulting sign is
$(-1)^{b(b-1)/2}$, proving \eqref{ar:det}.
The block decomposition and triangular transform preceding the theorem
prove \eqref{ar:spans} and \eqref{ar:complete}.

\subsection{Explicit reconstruction without matrix inversion}

The same positive recurrence supplies an exact reconstruction algorithm.
Given $C_b,\ldots,C_{2b-1}$, define
\[
            Y_r=\frac{2(-1)^{r-1}}{(2r-1)!}C_{b+r-1},
                       \qquad 1\le r\le b.
\]
Put $\eta_1=0$ and compute, for $1\le r<b$,
\begin{equation}\label{ar:eta}
       \eta_{r+1}=\frac{Y_r+\sum_{s=1}^rw_{rs}\eta_s}{d_{r+1}}.
\end{equation}
Then
\begin{equation}\label{ar:reconstruct}
 \begin{split}
 r_1&=-\frac{Y_b+\sum_{s=1}^bw_{bs}\eta_s}
                  {d_{b+1}\rho_{b+1}},\\
 r_s&=\rho_sr_1+\eta_s\quad(1\le s\le b),\\
 Z_k&=(-1)^{b-k}(2b-2k)!\,r_{b-k}\quad(0\le k<b).
 \end{split}
\end{equation}
The equations follow by writing the solution of
\eqref{ar:coefficient} as $r_s=\rho_sr_1+\eta_s$, then using
the final row. The denominator in \eqref{ar:reconstruct} is strictly
positive by the theorem. Equations \eqref{ar:backward} and
\eqref{ar:reconstruct} express every $Z_k$ as an explicit rational
linear combination of the first $q$ moments. Substitution into
\eqref{ar:block} and \eqref{ar:forward} recovers every later moment.

For illustration, the first five late-block determinants are
\[
 \begin{array}{c|ccccc}
 (p,q)&(2,4)&(4,6)&(6,8)&(8,10)&(10,12)\\\hline
 \det M_{p,q}&-11/48&-113/384&2509/5120&309713/307200&-7097067/2867200.
 \end{array}
\]
Their nontrivial numerators now have an exact origin in the positive
recurrence \eqref{ar:rho}. Finite rational checks verify these examples,
the general determinant formula through the supplied test range, the
rotation factors, and the reconstruction algorithm. The infinite proof
is the coefficient argument above, not the finite check.

\subsection{An infinite family of ordinary rational evaluations}

The reconstruction has an arithmetic consequence at every adjacent pair,
not just at the example printed below. Let $q=2b$, $p=q-2$, and denote
the first row of $M_{p,q}^{-1}$ by $(\lambda_b,\ldots,\lambda_{q-1})$.
These rational coefficients are explicitly computed by
\eqref{ar:eta}--\eqref{ar:reconstruct}. Define
\begin{equation}\label{ar:all-weights}
 w_r=\sum_{j=\max\{b,r\}}^{q-1}
     \lambda_j\binom{2j+1}{2r}2^{-2(j-r)},\qquad
 P_b(x)=\sum_{r=0}^{q-1}w_rx^{2r}.
\end{equation}
Put
\[
 u_0=\frac1{(p-1)(q-1)},\qquad
 u_1=-\frac1{24}\left(\frac q{p-1}+\frac p{q-1}\right).
\]

\begin{corollary}[Rational evaluations from every adjacent pair]
\label{ar:all-rational}
For every even $q\ge4$, with $p=q-2$ and $x_n=n+1/2$,
\begin{equation}\label{ar:all-ordinary}
 \begin{split}
 \sum_{n=0}^{\infty}\bigg\{
 &\frac{P_b(x_n)}{(p-1)!(q-1)!}
       \psi^{(p-1)}(n+1)\psi^{(q-1)}(n+1)\\
 &-u_0w_{q-1}x_n^2
      -u_0w_{q-2}-u_1w_{q-1}\bigg\}=-\frac12.
 \end{split}
\end{equation}
The series is absolutely convergent. Its polynomial has exact degree
$2q-2$, since
\begin{equation}\label{ar:top-weight}
       w_{q-1}
       =\frac{b(4b-1)\rho_b}{(2b+1)^2\rho_{b+1}}>0.
\end{equation}
\end{corollary}

\begin{proof}
By the defining inverse row, $\sum_{j=b}^{q-1}\lambda_jC_j=Z_0=-1/2$.
The inverse Bernoulli transform \eqref{ar:backward} then gives
$\sum_{r=0}^{q-1}w_rD_r=-1/2$.
For centered $x=n+1/2$, the two leading terms of the tail product are
\[
                 A_p(n)A_q(n)=u_0x^{2-p-q}+u_1x^{-p-q}
                                  +O(x^{-p-q-2}).
\]
Because $p+q=2q-2$, multiplication by $P_b$ leaves precisely the
quadratic and constant terms subtracted in \eqref{ar:all-ordinary}.
The remainder is $O(x^{-2})$. Both subtracted terms have continued
centered sum zero, so the convergent series equals the specified moment
combination. The polygamma normalization follows from the two even
orders. Finally the coefficient of $C_{q-1}$ in the reconstruction
of $Z_0$ is $b\rho_b/((2b+1)^2\rho_{b+1})$, as follows directly
from the last row of \eqref{ar:reconstruct}. Multiplication by
$\binom{2q-1}{2q-2}=2q-1=4b-1$ proves \eqref{ar:top-weight}.
\end{proof}

\subsection{A convergent rational polygamma identity at orders four and six}

The first nontrivial new adjacent example after $(2,4)$ is $(4,6)$.
In this case \eqref{ar:entries} and the exact reconstruction give
\begin{equation}\label{ar:46matrix}
 \begin{pmatrix}C_3\\C_4\\C_5\end{pmatrix}
 =\begin{pmatrix}
 0&3/4&1/4\\5/12&2/3&-1/12\\1&-3/8&1/12
 \end{pmatrix}
 \begin{pmatrix}-1/2\\\zeta(2)\\\zeta(4)\end{pmatrix},
\end{equation}
and the inverse matrix is
\[
 \frac1{339}
 \begin{pmatrix}-28&180&264\\136&288&-120\\948&-864&360\end{pmatrix}.
\]
The first row, followed by \eqref{ar:backward}, yields
\begin{equation}\label{ar:46moment}
 67D_0+15516D_1+324320D_2+1434496D_3
       +1601280D_4+371712D_5=-21696.
\end{equation}
This has an ordinary, absolutely convergent realization. Put
\[
 \begin{split}
 P(x)={}&67+15516x^2+324320x^4+1434496x^6\\
        &+1601280x^8+371712x^{10},\qquad x_n=n+\tfrac12.
 \end{split}
\]
Then
\begin{equation}\label{ar:46ordinary}
 \boxed{\displaystyle
 \sum_{n=0}^{\infty}\left\{
  \frac{P(x_n)}{720}\psi^{(3)}(n+1)\psi^{(5)}(n+1)
       -\frac{123904}{5}x_n^2-\frac{316928}{5}\right\}=-21696.}
\end{equation}
To check the subtraction and the meaning of this identity, the centered
Hurwitz expansion gives
\[
 A_4(n)A_6(n)
   =\frac1{15}x_n^{-8}-\frac7{60}x_n^{-10}+O(x_n^{-12}).
\]
Multiplication by $P$ leaves precisely the two nonnegative powers
displayed in \eqref{ar:46ordinary}; the remaining bracket is
$O(x_n^{-2})$. Their continued centered sums are zero because
$\zeta(-2,1/2)=\zeta(0,1/2)=0$. Thus the absolutely convergent
sum equals the meromorphically continued moment combination in
\eqref{ar:46moment}. Finally
$A_4(n)=\psi^{(3)}(n+1)/6$ and
$A_6(n)=\psi^{(5)}(n+1)/120$, explaining the factor $720$.
Every summand is an ordinary real number; no finite-part instruction
remains in \eqref{ar:46ordinary}.

There is also a finite telescoping proof of \eqref{ar:46ordinary}
which is independent of the continued-moment formulas. Define
\begin{align*}
 K(N)&=-3584N^7+23040N^9+33792N^{11},\\
 P_4(N)&=128N(N+1)
       (44N^4+88N^3+67N^2+23N-14),\\
 P_6(N)&=128N^3(N+1)^3(33N^2+33N-14),\\
 R_0(N)&=-\frac{74624}{15}N^3+4928N^2-\frac{244288}{15}N.
\end{align*}
Put
\[
 \mathcal F(N)=K(N)A_4(N)A_6(N)
      +P_4(N)A_4(N)+P_6(N)A_6(N)+R_0(N).
\]
The identities $A_p(N-1)=A_p(N)+N^{-p}$, together with direct
polynomial multiplication, give
\begin{equation}\label{ar:46telescoping}
 \mathcal F(N)-\mathcal F(N-1)
  =P(N-\tfrac12)A_4(N-1)A_6(N-1)
      -\frac{123904}{5}(N-\tfrac12)^2-\frac{316928}{5}.
\end{equation}
This equality is a rational-function identity in $N$ and in two
indeterminates replacing $A_4(N)$ and $A_6(N)$; no zeta relation is
used to verify it. Since $\mathcal F(0)=0$, the partial sum of
\eqref{ar:46ordinary} through $n=N-1$ is exactly $\mathcal F(N)$.
The standard uncentered Hurwitz expansions
\[
 \begin{split}
 A_4(N)&=\frac1{3N^3}-\frac1{2N^4}+\frac1{3N^5}+O(N^{-7}),\\
 A_6(N)&=\frac1{5N^5}-\frac1{2N^6}+\frac1{2N^7}+O(N^{-9})
 \end{split}
\]
give $\mathcal F(N)=-21696+O(N^{-1})$. All positive powers
cancel; the constant contributions from $KA_4A_6$, $P_4A_4$, and
$P_6A_6$ are respectively $-20224$, $-448$, and $-1024$.
This proves the evaluation by an independent finite identity followed
by an elementary asymptotic limit.

\paragraph{Scope.}
This proves the source conjecture for two infinite classes: the source's
diagonal class and the new adjacent-order class. The phrase ``optimal
prefix'' refers to the formal coordinate presentation; it does not
prove that $q$ actual numerical constants are independent. For larger
gaps the rotated equation has several free initial coefficients.
Positivity of its scalar recurrence alone does not establish
nonsingularity of the resulting multi-column boundary matrix. That
remaining distinction is essential and prevents an unsupported
extension of the present theorem to all unequal pairs.

% Optional continuation after adjacent_tail_prefix.tex. All labels ar:.
\subsection{A smaller boundary determinant for a fixed gap}
\label{ar:general-gap}

The rotation argument has a useful exact extension even though it does
not settle the complete unequal-order conjecture. Write
$q=2b$, $p=2(b-c)$, with $1\le c<b$, so the differential order in
the source operator is $m=2c$. For integers $s\ge c+1$ put
\[
 A_s=\frac{(2s)!}{(2s-2c)!},\qquad
 B_s=\frac{(2s-1)!}{(2s-2c-1)!},\qquad
 d_s=2(A_s+B_s),
\]
and put $A_s=0$ for $1\le s\le c$. Define $c$ fundamental rational
sequences by
\begin{align}
 R_{s,j}&=\delta_{sj},&&1\le s,j\le c,\notag\\
 R_{s,j}&=\frac1{d_s}\sum_{u=1}^{s-1}
       \alpha_{s-u}(B_s+A_u)R_{u,j},
       &&s\ge c+1,\quad 1\le j\le c.
       \label{ar:fundamental}
\end{align}
Every generated entry $R_{s,j}$ is strictly positive. The relevant
boundary matrix has size $c$, independent of $b$:
\[
             F_{b,c}=(R_{b+i,j})_{1\le i,j\le c}.
\]

\begin{proposition}[Fixed-gap boundary reduction]\label{ar:gap-det}
In the normalization \eqref{ar:M},
\begin{equation}\label{ar:gap-factor}
 \det M_{2(b-c),2b}
  =\frac{(-1)^{b(b-1)/2}}{4^bb!}
       \left(\prod_{s=c+1}^{b+c}d_s\right)\det F_{b,c}.
\end{equation}
In particular the optimal-prefix conjecture for a given gap $2c$
is equivalent to $\det F_{b,c}\ne0$ for every integer $b>c$.
When $c=1$, $R_{s,1}=\rho_s$ and this is
Theorem~\ref{ar:main}.
\end{proposition}

\begin{proof}
For $R(x)=Q(ix)=\sum r_sx^{2s}$, rotate the source operator to
\[
 \bigl(R^{(2c)}-R^{(2c)}(0)\bigr)\cot(x/2)
        +(R\cot(x/2))^{(2c)}.
\]
The coefficient of $x^{2k-1}$, with $s=k+c$, is
\begin{equation}\label{ar:general-coefficient}
       d_sr_s-\sum_{u=1}^{s-1}\alpha_{s-u}(B_s+A_u)r_u.
\end{equation}
The pole $2/x$ gives $d_s$; every other cotangent coefficient has
the negative sign in \eqref{ar:cot}. The constant derivative
subtraction again makes $A_c=0$, and all lower derivatives already
vanish. This proves the displayed coefficient formula.

The first $b-c$ vanishing rows successively determine
$r_{c+1},\ldots,r_b$ from $r_1,\ldots,r_c$ by
\eqref{ar:fundamental}. The remaining $c$ rows, with
$r_{b+1}=\cdots=r_{b+c}=0$, are equivalent to the vanishing of
the boundary coordinates $F_{b,c}(r_1,\ldots,r_c)^{\mathsf T}$.
For the determinant, the first elimination contributes
$\prod_{s=c+1}^bd_s$. The remaining boundary block differs from
$-\operatorname{diag}(d_{b+1},\ldots,d_{b+c})F_{b,c}$ by a
unit lower triangular row transformation; its determinant is
$(-1)^c\prod_{s=b+1}^{b+c}d_s\det F_{b,c}$.
Moving the $c$ initial columns past the other $b-c$ columns contributes
$(-1)^{c(b-c)}$.

Finally, the rotation gives the row factor
$2(-1)^{r-c}/(2r-1)!$ and the same reversed column factors
$(-1)^s/(2s)!$ as before. The resulting conversion from the
coefficient matrix in \eqref{ar:general-coefficient} to $M$ has
factor $(-1)^{b(b-1)/2+bc}/(4^bb!)$. All additional signs multiply
to $(-1)^{c(1-c)}=1$, proving \eqref{ar:gap-factor}.
\end{proof}

The positivity of all individual entries in \eqref{ar:fundamental}
does not prove the determinant nonzero. There is not even a uniform
positive determinant sign for every fixed gap: exact evaluation at
$c=3$ gives $\det F_{6,3}<0$ and $\det F_{7,3}>0$. Both are nonzero.
The remaining research question is therefore a genuine boundary
determinant problem, rather than a direct extension of the one-dimensional
positivity proof.

% END sections/08_adjacent_tails.tex


% BEGIN sections/09_audit_and_questions.tex
\section{Source audit, proof status, and further research}
\label{audit:section}

\subsection{What has changed relative to the source questions}

The baseline for this continuation is the repository revision
\[
\texttt{\sourcecommit},
\]
dated 11 October 2026. The two paths named
in the research request were inspected at that revision. Six current
incoming archives were extracted for comparison, and relevant canonical
chapters, the manuscript overview, and its proof-status records were
consulted. The source manifest in the package records their paths and
hashes. This fixes what \emph{existing} means in the comparison; it does
not imply that all historical drafts were audited.

The main advancements have different logical forms and should be
integrated accordingly.
\begin{enumerate}
\item The higher-power request for differential-polynomial finite parts
receives an all-polynomial theorem, Theorems~\ref{poly:ML},
\ref{poly:filter}, and \ref{poly:bridge}. The new scalar projection
calculation identifies exactly which polynomial part produces the
nondecaying residue and its Stieltjes correction. Mixed monomials have
a finite, individually convergent word reduction by
Theorem~\ref{poly:words}. The scalar power case is recovered prior work.

\item The canonical report's question about integration by parts at all
weights has the exact direct-sum solution in
Theorem~\ref{poly:IBP}. The endpoint map is retained in
\eqref{poly:IBPintegral}. This solves the formal quotient problem;
classification of the numerical residue-coordinate image still requires
relations among the resulting convergent spectral derivatives.

\item The scalar Gamma--Appell question is sharpened and partly settled
at every degree by Theorem~\ref{ac:main}: the highest height-one Laurent
coefficient is the only \emph{singular} coordinate needed. The explicit
sixth and seventh formulas implement the requested higher-degree
calculations. Elimination of all convergent mixed coordinates is a
different assertion, and remains open.

\item The common-shift problem receives normalized primitives and a
proved endpoint limit in Theorem~\ref{poly:primitivethm}. This establishes
the precise subtraction and its limiting constant. A reduction of every
shifted primitive to a chosen multiple-Hurwitz or Barnes basis is not
proved here.

\item Q8 of the continuous-depth report receives the requested all-order
Gamma--Stieltjes--$Z_\rho$-jet answer in
Theorem~\ref{lj:thm:oddjets}. The reflected jets are supplied with their
own convergent evaluator in Theorem~\ref{lj:thm:Euler}, and the result
extends to ordinary Gamma-weighted polylogarithm derivatives.

\item Conjecture \texttt{conj:firstq} of the canonical fifth-moment
report is proved for all $p=q-2$ by Theorem~\ref{ar:main}.
The general unequal-order conjecture is reduced to an explicit
fixed-gap boundary determinant by Proposition~\ref{ar:gap-det}, but
its nonvanishing is not proved for all larger gaps. The rational
polygamma-product evaluations in Section~\ref{ar:section} include
an independently telescoping example.
\end{enumerate}

\subsection{Corrections and editorial recommendations}

No newly identified false theorem is asserted for the inspected portions
of the six incoming reports. The existing cautions about finite parts,
formal rank, and differentiation of discrete identities are sound.
The following changes are proposed for their integration and for the
research-question lists.

First, update the optimal-prefix conjecture to record the proved
neighboring-order family while retaining the arbitrary-gap statement as a
conjecture. Its determinant and the exact binomial omission convention
should be copied together: retaining the lower-index-zero summand in
\eqref{ar:entries} changes the matrix and invalidates the proof.
Do not replace formal rank by independence of the actual tail
moments or of the even zeta values.

Second, state integration-by-parts reductions as
$\cI(P)=E(Q)+\cI(N)$, with $Q$ at the prescribed weight. For example,
$\cI(u_1^2+u_0u_2)=-\zeta(3)$, not zero. This is an explicit
counterexample to discarding the endpoint constant; the source question
already correctly requires that constant. The new normal-form theorem
provides the requested implementation of that requirement.

Third, descriptions of Gamma--Appell cancellation should distinguish the
singular coordinates from convergent ones. The formulas
\eqref{ac:D6} and \eqref{ac:D7} contain genuine free-order derivatives
of admissible multiple zeta functions. Neither their appearance nor a
failed numerical relation search proves that they are irreducible.
Conversely, calling the singular support a singleton does not state
that the whole answer contains only one nonclassical constant.

Fourth, add the local identity \eqref{lj:eq:local} next to the source's
odd-lattice formula before discussing order derivatives. The known
polynomial values in the discrete index determine the $q=0$ case only.
The reflected derivatives $\partial_s^jZ_\rho(-2m)$ are required for
$q>0$. The same care applies to the normalization factor $2^{1-s}$
in \eqref{lj:eq:Imaster}: its derivatives produce the explicit
$\log 2$ terms in the Stieltjes examples.

Finally, retain the manuscript's existing distinction between its
rejected old Gaussian $S_8$ vector and its revised $S_8$ candidate.
The Gaussian $S_6$ and revised $S_8$ remain conjectural
\cite{Repo}. The symbols $\mathsf S_6,\mathsf S_7$ in this article
are Gamma--Appell finite-part moments, unrelated to those Gaussian
ladder labels. None of the present results changes the Gaussian
proof status.

\subsection{Verification and the role of the accompanying artifacts}
\label{audit:verification}

The proofs in the article establish the all-order statements. The
package also supplies exact finite algebra and independent numerical
diagnostics. Their distinct purposes are recorded in machine-readable
result files.

The adjacent-tail suite checks the source binomial matrix against the
rotated cotangent operator, the positive recurrence against exact
determinants, the full moment reconstruction, the fixed-gap
factorization, and the explicit telescoping certificate.
It records 2604 exact rational assertions and 90 independent symbolic
cotangent-derivative assertions, with adjacent orders through $q=36$.
A deliberately altered coefficient is rejected. Finite examples do not
prove the all-order determinant statement; the positive recurrence does.

The Gamma--Appell suite constructs finite Laurent data, the tail defects,
and the strict word maps. Its weight-six certificate contains 49
shuffle, stuffle, and duality equations on 16 admissible words, together
with an invertible rational subsystem. The separate replay verifies the
certificate and the free-order directional reductions. Integer-index
duality is used for ordinary values only. An additional independent
calculation using ordered set partitions agrees with both printed
directional formulas.

The Lerch suite compares exact Bell coefficients with direct logarithmic
series, checks finite Euler values and trivial zeros, and compares
negative-order Euler sums with positive-order Gamma-moment integrals.
It records 48 exact cumulant checks, 45 finite Euler evaluations
(including 20 odd trivial zeros), two differentiated-product checks,
30 independent mixed-jet numerical comparisons, and direct Stieltjes
derivative checks. The latter comparisons include noninteger and complex
depth. Their largest recorded relative residual is approximately
$1.17\times10^{-71}$.

The differential-polynomial suite verifies all 270 formal monomial
decompositions of weights 2--12, with 22 two-sided exact matrix inverse
identities. It checks 128 mixed and 10 scalar Laurent residues through
weight ten; 48 bilinear pairs, including every applicable zero-order
case; 169 finite harmonic identities; and explicit endpoint constants.
Its numerical companion computes finite parts using a convergent Taylor
integral near zero and ordinary quadrature away from zero. It compares
two splitting points, tests nonlinear endpoint identities, checks
normalized primitives at real and complex arguments, and differentiates
ordinary common-shift integrals.

All numerical records use arbitrary-precision floating-point arithmetic;
they are not interval certificates. The article does not claim
proof-assistant formalization. The symbolic certificates are exact
finite objects, but the analytic continuation, convergence, and
normalization lemmas remain mathematical proofs in the text.

\subsection{Further research questions}
\label{audit:questions}

The following questions seek identities, finite classifications, and
functional relations. Estimates would serve to justify operations or
certificates, rather than being the principal research target.

\subsubsection{The boundary determinant for arbitrary even gaps}

For each fixed $c\ge2$, prove or disprove
$\det F_{b,c}\ne0$ for every $b>c$, with the explicit fundamental
sequences in \eqref{ar:fundamental}. The case $c=2$ is the first
unresolved family. The reduction is useful because the determinant size
depends on the gap, not the upper tail order. A successful proof could
come from a discrete Wronskian identity, a recurrence for the minors,
or a factorization in a suitable polynomial basis. A uniform positivity
argument for all gaps is unavailable: the sign change already recorded
at $c=3$ rules out that simplest claim.

\subsubsection{Systematic rational telescoping for tail products}

The explicit $(4,6)$ identity has a certificate involving only a
polynomial multiple of $A_4(N)A_6(N)$, two single-tail terms, and a
polynomial. Determine whether every adjacent-order reconstruction
identity admits a certificate of this form, and derive its coefficients
directly from the recurrence for $\rho_s$. Such a theorem would replace
spectral continuation by a finite telescoping proof for the whole
ordinary-series family. A useful first output is a bound on the required
polynomial degrees proved from the shift recurrence for the tails.

\subsubsection{The spectral image of the canonical differential basis}

Theorem~\ref{poly:IBP} classifies differential monomials before arithmetic
relations are imposed. Apply the convergent-word map to this basis and
determine its kernel modulo a precisely stated family of free-order
stuffle relations and differential identities. The first interesting
weights are those containing several nonlinear mixed basis elements,
such as $u_0^2u_1^2$ and $u_1^3$ at weight six. The desired result is a
finite presentation with a proved rank in that presentation. Arithmetic
independence of its evaluated numbers should remain a separate question.

\subsubsection{Relations among the sixth and seventh mixed directions}

Determine which combinations in \eqref{ac:D6} and \eqref{ac:D7}
reduce to lower depth, ordinary zeta derivatives, or named
height-one Laurent coefficients. Symmetric free-order stuffle relations
already explain part of the reduction; the surviving antisymmetric
directions require additional identities. Tornheim functional relations,
parameter-dependent iterated integrals, and deformation of known
double-zeta identities are concrete candidates. An identity known only
at integer orders cannot be differentiated for this purpose.

\subsubsection{Higher spectral derivatives of the residue calculus}

The present finite-part map selects the coefficient of $s$ in a
Dirichlet germ. Introduce logarithmically weighted spectral functionals
whose ordinary representatives select $[s^q]$ for $q\ge2$.
Determine whether the same scalar projection controls all their
Stieltjes corrections, and identify a canonical family of normalized
spatial kernels realizing them. For mixed monomials the strict
upper--middle--lower construction already gives convergent higher
partial derivatives; the remaining issue is the exact spatial
normalization and its endpoint constants.

\subsubsection{Rational shifts and multiple-Hurwitz coordinates}

For $I_P(a)$ and $\mathcal P_P(a)$ at rational $a>0$, find finite
formulas in a specified algebra of multiple-Hurwitz derivatives and
cyclotomic polylogarithms. The normalized kernels in
\eqref{poly:primitivekernel} and the endpoint law
\eqref{poly:shiftlimit} fix the constants that such a formula must
respect. Reflection and multiplication combinations may simplify before
individual values do. A complete result should also explain the relation
to the Gamma-quotient normalization with a general base point.

\subsubsection{Additional reductions at rational noninteger depth}

For fixed rational noninteger $\rho>0$, investigate the reflected
coordinates $\partial_\rho^r\partial_s^qZ_\rho(-2m)$.
Theorem~\ref{lj:thm:Euler} already evaluates them by normally convergent
series of finite harmonic-logarithmic expressions. The further
arithmetic question is whether selected depths admit finite expressions
in a predeclared collection of Hurwitz, Gamma, or cyclotomic constants.
Integer-depth termination at $q=r=0$ does not persist automatically
after differentiation. Identifying exactly when an additional finite
reduction exists would be a substantive strengthening of the current
lattice theorem.

\subsubsection{Independent twists and noncentered Gamma weights}

Extend the local lattice construction to the noncentered and tilted
families while keeping order, shift, and depth as independent complex
variables. The centered theorem uses the even Taylor structure and a
single sine factor. Away from the center, the reflection system can
couple two branches or two transforms. Derive the corresponding local
matrix identity and determine its removable values before extracting
mixed jets. The desired output is a finite formula for ordinary
Gamma-weighted root-of-unity polylogarithm derivatives, with every
branch and normalization specified.

\subsubsection{Four or more harmonic tails}

For products of an even number of even-order tails, develop a finite
coordinate system and prove an initial-moment generation theorem.
Centered parity still controls which spectral values can be singular,
but the two-tail rank calculation does not establish the higher-product
claim. The relevant differential operators should contain products of
cotangent kernels and higher pole multiplicities. Exact finite
reconstruction and ordinary telescoping identities would be natural
outputs. Odd numbers of tails require a separate finite-part
normalization because their centered parity changes.

\subsubsection{Functional certificates for the remaining ladders}

Continue the exact functional-equation search for the manuscript's
Gaussian $S_6$ and revised $S_8$ candidates, keeping the evaluated
identity distinct from the formal symbol calculation. The new
convergent spectral coordinates suggest additional Mellin and
polylogarithmic representations, but no bridge to those particular
ladder vectors is proved here. A successful next step should supply an
explicit functional certificate, a proved obstruction in a stated
relation system, or a rigorously nonzero evaluated residual.

\subsubsection{Formal certificates with a small analytic interface}

The finite algorithms admit a useful separation: exact Laurent
extraction and word multiplication, rational matrix reconstruction,
and the analytic lemmas for normal convergence and finite parts.
Formalize the first layer independently, then connect it through proved
interfaces to the scalar residue filter, the normalized
Mittag--Leffler theorem, and the Euler holomorphy theorem.
The purpose is a reusable proof pipeline for further identities,
with each numerical diagnostic remaining a separately identified
experiment.

% END sections/09_audit_and_questions.tex


% BEGIN references.tex
\begingroup
\small
\begin{thebibliography}{99}
\addcontentsline{toc}{section}{References}

\bibitem{Repo}
ProveIt Contributors,
\emph{Polylogarithms and their Arithmetic Bridges},
canonical manuscript and editorial records, revision
\texttt{\sourcecommit}, 11 October 2026.
\url{https://github.com/VladimirReshetnikov/ProveIt/tree/f7e35be357654198c1546fd40f363694c5bcc2cc/Analysis/Polylogarithms/docs/manuscript}.

\bibitem{Higher}
ProveIt research continuation,
\emph{Higher Digamma Powers and Ordered Gamma Sums}, 11 October 2026.
Incoming archive at the same revision; especially the all-power
Mittag--Leffler theorem, the finite harmonic-word algorithm, and the
differential-polynomial, common-shift, and higher-moment questions.
\url{https://github.com/VladimirReshetnikov/ProveIt/blob/f7e35be357654198c1546fd40f363694c5bcc2cc/docs/incoming/ProveIt_Higher_Digamma_Powers_2026-10-11.zip}.

\bibitem{Fifth}
ProveIt research continuation,
\emph{Canonical Digamma Polynomials and the Fifth Harmonic Bridge:
Laurent Coefficients, Polygamma Primitives, and Finite Reconstruction
of Even-Tail Moments}, 11 October 2026.
Incoming archive at the same revision; Sections 4, 7, and 9,
including Conjecture \texttt{conj:firstq}.
\url{https://github.com/VladimirReshetnikov/ProveIt/blob/f7e35be357654198c1546fd40f363694c5bcc2cc/docs/incoming/ProveIt_Canonical_Digamma_Fifth_Moment_2026-10-11.zip}.

\bibitem{Mixed}
ProveIt research continuation,
\emph{Mixed Resolvent Products, Higher Digamma Powers, and Stieltjes
Collisions: Exact Identities and Normalized Finite Parts},
11 October 2026. Incoming archive at the same revision.
\url{https://github.com/VladimirReshetnikov/ProveIt/blob/f7e35be357654198c1546fd40f363694c5bcc2cc/docs/incoming/ProveIt_Mixed_Resolvents_Harmonic_Powers_2026-10-11.zip}.

\bibitem{Continuous}
ProveIt research continuation,
\emph{Continuous-Depth Lerch Calculus: Gamma Kernels, All-Order Harmonic
Identities, Stieltjes Convolutions, and Anchored Primitives},
10 October 2026. Incoming archive at the same revision; the centered
reflection theorem, Gamma convolution, and final research question Q8.
\url{https://github.com/VladimirReshetnikov/ProveIt/blob/f7e35be357654198c1546fd40f363694c5bcc2cc/docs/incoming/ProveIt_Continuous_Depth_Lerch_2026-10-10.zip}.

\bibitem{Tilted}
ProveIt research continuation,
\emph{Tilted Lerch Convolutions and Binary Harmonic Cancellation:
Confluent Zeta Identities, Independent Spectral Jets, Gamma Primitives,
and Mixed-Parity Exact Evaluations}, 10 October 2026.
Incoming archive at the same revision.
\url{https://github.com/VladimirReshetnikov/ProveIt/blob/f7e35be357654198c1546fd40f363694c5bcc2cc/docs/incoming/ProveIt_Tilted_Lerch_Binary_Cancellation_2026-10-10.zip}.

\bibitem{Algebraic}
ProveIt research continuation,
\emph{Algebraic Gauss Residues and Centered Harmonic Reflection},
incoming archive at the same revision.
\url{https://github.com/VladimirReshetnikov/ProveIt/blob/f7e35be357654198c1546fd40f363694c5bcc2cc/docs/incoming/Algebraic_Gauss_Centered_Harmonic_Reflection.zip}.

\bibitem{DLMFGamma}
NIST Digital Library of Mathematical Functions,
Chapter 5, especially \S\S5.7 and 5.15,
Gamma expansions, polygamma recurrences, and special values.
\url{https://dlmf.nist.gov/5.7};
\url{https://dlmf.nist.gov/5.15}.

\bibitem{DLMFZeta}
NIST Digital Library of Mathematical Functions,
\S25.11, Hurwitz zeta function.
\url{https://dlmf.nist.gov/25.11}.

\bibitem{DLMFLerch}
NIST Digital Library of Mathematical Functions,
\S25.14, Lerch's transcendent.
\url{https://dlmf.nist.gov/25.14}.

\bibitem{DLMFEuler}
NIST Digital Library of Mathematical Functions,
\S3.9(ii), Euler's transformation of series.
\url{https://dlmf.nist.gov/3.9#ii}.

\bibitem{IKZ}
K. Ihara, M. Kaneko, and D. Zagier,
Derivation and double shuffle relations for multiple zeta values,
\emph{Compositio Mathematica} \textbf{142} (2006), 307--338.
\url{https://doi.org/10.1112/S0010437X0500182X}.
Author-hosted full text:
\url{https://www2.math.kyushu-u.ac.jp/~mkaneko/papers/34Derivation_double_shuffle.pdf}.

\bibitem{RamanujanFourier}
D. Chakrabarti and G. K. Srinivasan,
\emph{On a remarkable formula of Ramanujan}, 2012.
\url{https://arxiv.org/abs/1206.5122}.
The Gamma-product Fourier normalization used here is also derived
directly from the beta integral in Section~\ref{lj:sec:main}.

\end{thebibliography}
\endgroup

% END references.tex

\end{document}
